\begin{theorem}[Source result P3K.118 under continuation review]\label{night:P3K.118}
\StatementRecord{night:P3K.118}{open}{}{the original source, point line, roots, characters, reference, selected fan and scalar range stated here}
Retain the original marked source, selected whole nodal fibre and point line
\[
y^2=x^3+tx+t^2,\qquad G=(0,t),\qquad
q=t+4/27,\qquad \mathcal M_a=\mathcal O([aG]-[O]),\qquad ab=6,
\]
with source translation by \(-bG\). Its two original additive caps,
including their full Cartier maps, finite characters and prescribed
annular references, have orders in the following four domains:
\[
\begin{array}{c|c|c}
(a,b)&m=m_\infty&h=m_0\\ \hline
(1,6)&3l,\ l\ge1&4k,\ k\ge1\\
(2,3)&3l,\ l\ge1&4k+1,\ k\ge0\\
(3,2)&3l+2,\ l\ge0&4k,\ k\ge1\\
(6,1)&3l+2,\ l\ge0&4k+1,\ k\ge0.
\end{array}
\]
The prescribed norms satisfy \(hn_\infty+mn_0=\varepsilon\), with
\(\varepsilon=\pm1\). The additional ordinary total norm is zero.
Retain every prescribed selected integral unimodular fan and the
original strictly contracting scalar domain. These objects have the
original six-sphere topology through the full F46 source/cap and
common-blowup maps.

The holomorphic automorphism group of each such marked object is
exactly its faithful physical \(\mathbb C^*\). Every holomorphic
scalar endpoint fixes the original base \(t\), has positive affine
linear part on the regular two-tori, and has equal original scalar
values at its two ends. Its possible integer scalar windings are
exactly
\[
J_0\mathbb Z,\qquad J_0=\frac{mh}{b}.
\]
At a fixed winding the full holomorphic endpoints differ precisely
by the extending physical scalar. This classification includes both
root caps and their complete reference words.

Retain the rational positive-\(G\) line operator \(U\), the actual
reference constants \(C_a,c_0\), and the full original identities
\[
B_0=C_a\frac q{t^3}U^{-b},\qquad
h_0=c_0q^\varepsilon t^{-mn_0}.
\]
On \(\lambda=r^{mh}\), the primitive forward endpoint is
\[
E_r=\gamma_0r^{-\varepsilon J_0}t^{d_0}U^\varepsilon,
\qquad
d_0=\frac{3\varepsilon-mn_0}{b},\qquad
\gamma_0^b=c_0C_a^{-\varepsilon},
\qquad r\longmapsto\zeta_b r.
\]
Its inverse is evaluated at the actual inverse source point. Its
\(b\) chronological copies give exactly
\(h_0B_\lambda^{-\varepsilon}\), with all original constants retained.
The complete K67 parameter-family comparison is the physical map
\[
F_r=r^{\varepsilon J_0(b-1)/2},\qquad
F_{\zeta_b r}E_r=
\exp\!\left(\frac{\pi i\varepsilon J_0(b-1)}b\right)E_rF_r.
\]
The exponent is integral on every displayed domain. The original
family is the literal degree-\(b\) pullback of the completed shorter
family over \(w=r^b\), \(\lambda=w^{J_0}\), including both caps,
the selected Cox maps, physical line, stable tangent and all seams.

For one positive primitive scalar circle, the initial physical pair
\((k_{\rm phys},\ell_{\rm phys})\) and index \(S_0\) are
\[
\begin{array}{c|c|c|c}
(a,b)&k_{\rm phys}&\ell_{\rm phys}&S_0\\ \hline
(1,6)&kl&4kl&0\\
(2,3)&2l(4k+1)&8l(k+1)&kl\\
(3,2)&9k(3l+2)&12k(l+2)&k(4l+2)\\
(6,1)&18(3l+2)(4k+1)&24(kl+2k+l+1)&35kl+22k+8l+5.
\end{array}
\]
The actual held-quaternion native pair is
\[
(\mu,\nu)=(\ell_{\rm phys}/4,\,
k_{\rm phys}-\ell_{\rm phys}/4).
\]
For the vertical fan \(T_j\), the full primitive data are
\[
S_j=S_0-\frac{J_0}{2}j(j-1),\qquad
k_j=k_{\rm phys}-3J_0j(j-1),\qquad
\ell_j=\ell_{\rm phys}.
\]
For every selected lattice-basis fan \(P\), the complete terminating
algorithm RJC.78--80, equivalently the full native maps and
nine-variable target systems JFP.41--43, gives its actual integer
return \(S(P)\). This includes every full source coefficient, original
height, ordinary product, physical normal and cap reset.

Every holomorphic scalar path of winding \(sJ_0\) and every
holomorphic closing at that winding has return
\(sS(P)\pmod{28}\) on its prescribed fan. Conversely every
integer \(s\) is attained by the explicit shorter-family annular
pullback. The exact images over all selected fans and all such
paths are
\[
\begin{array}{c|c|c}
(a,b)&\text{holomorphic winding group}&
\text{complete all-fan return image}\\ \hline
(1,6)&2kl\mathbb Z&kl\mathbb Z/28\\
(2,3)&l(4k+1)\mathbb Z&l\mathbb Z/28\\
(3,2)&2k(3l+2)\mathbb Z&k\mathbb Z/28\\
(6,1)&(3l+2)(4k+1)\mathbb Z&\mathbb Z/28.
\end{array}
\]
For the positive seam \((x,1)\sim(\Phi(x),0)\), the oriented
mapping torus is
\[
(S^1\times S^6)\mathbin{\#}\Sigma_{sS(P)}.
\]
The inverse seam gives \(\Sigma_{-sS(P)}\). For each of the two
spin structures the corresponding Eells--Kuiper values are
\(sS(P)/28\) and \(-sS(P)/28\) in \(\mathbb Q/\mathbb Z\).
The two signed sphere equations retain the same positive scalar
parameter orientation throughout these statements.

All assertions retain the four stated root progressions, their
original norm representatives, every finite character, full annular
reference and selected fan. The scope is the fixed original source
and zero additional ordinary total norm.
\end{theorem}
\begin{proof}[Supplied source argument; continuation review in progress]

% Expanded verified cp36 input: notes/part3_e6_iii_joint_endpoints_proof.tex
\par\bigskip\noindent\textbf{Complete supporting proof 1: RJC.}\par


% Expanded verified cp36 input: notes/part3_e6_iii_joint_endpoints_rjc_proof.tex
\emph{Retained preparation history.} The primary source\textquotesingle s introductory and final preparation-stage audit sentences are reproduced verbatim. The complete independent derivation and geometric audit follow as JFP.1-\/-12. The exact current gates and their source locks are recorded in \path{records/k118_export_manifest.json}. The completed reciprocal full-source review is included in that gate.

\par

\medskip\noindent\textbf{The original E6 source with both additive caps: endpoint continuation}

\par\nobreak

This continuation retains K113-\/-115 and the full F46 joint-cap construction.
The original source is
\[
 E:\quad y^2=x^3+tx+t^2,\qquad G=(0,t),\qquad
 t_c=-4/27,\quad q=t+4/27.                         \tag{RJC.1}
\]
This is the full primary continuation on the four actual two-cap
domains below. It retains every point-line module, root comparison,
finite character and prescribed annular reference. It constructs the
primitive endpoint, classifies all holomorphic scalar endpoints, and
carries their parameter families to every lattice-basis fan. The
independent full audit and the manuscript statement remain pending.
The earlier derivation, certificate and output are preserved exactly in
\path{records/attempts/e6_iii_joint_before_complete_proof/}.

\par

\medskip\noindent\textbf{RJC.1. The actual four root domains and the signed sphere equation}

\par\nobreak

Write \(m=m_\infty\), \(h=m_0\), \(M=mh\). The actual root domains of
F46, including the finite root rather than its order-one special case, are
\[
\begin{array}{c|c|c|c}
 (a,b)&m&h&\text{minimal root fibres}\\ \hline
 (1,6)&3l,\ l\ge1&4k,\ k\ge1&\text{smooth, smooth}\\
 (2,3)&3l,\ l\ge1&4k+1,\ k\ge0&\text{smooth, III}\\
 (3,2)&3l+2,\ l\ge0&4k,\ k\ge1&\text{IV, smooth}\\
 (6,1)&3l+2,\ l\ge0&4k+1,\ k\ge0&\text{IV, III}.
\end{array}                                                        \tag{RJC.2}
\]
The two retained norms satisfy
\[
 hn_\infty+mn_0=\varepsilon,\qquad \varepsilon\in\{1,-1\}.             \tag{RJC.3}
\]
This is the actual two-cap sphere equation proved by all integral maps
F46.J.7-\/-12. It forces \(\gcd(m,h)=1\) and each local primitive norm.
The ordinary order-one framing in this first continuation has zero total
additional norm. Its full annular reference and both finite characters
are retained. Additional ordinary clutch slots have their separate
K32 exponents and will be accounted for explicitly before that scope is
included.

The minimal root substitutions are, at infinity,
\[
\begin{array}{c|c|c}
m&(t,x,y)&\text{root equation}\\ \hline
3l&(z^{-3l},z^{-2l}X,z^{-3l}Y)&Y^2=X^3+z^lX+1\\
3l+2&(z^{-(3l+2)},z^{-2(l+1)}X,z^{-3(l+1)}Y)
 &Y^2=X^3+z^{l+2}X+z^2,
\end{array}                                                        \tag{RJC.4}
\]
and at zero,
\[
\begin{array}{c|c|c}
h&(t,x,y)&\text{root equation}\\ \hline
4k&(z^{4k},z^{2k}X,z^{3k}Y)&Y^2=X^3+X+z^{2k}\\
4k+1&(z^{4k+1},z^{2k}X,z^{3k}Y)
 &Y^2=X^3+zX+z^{2k+2}.
\end{array}                                                        \tag{RJC.5}
\]
The good III central point is \(T=(0,0)\) on \(Y^2=X^3+X\).
The good IV-star central point is \((0,1)\) on \(Y^2=X^3+1\).
Their orders are respectively two and three, by the tangent divisors.
The inverse substitutions over \(z\ne0\) are multiplication by the
reciprocal displayed powers. The deck acts by multiplying \(z\) by
\(\zeta_{m_i}\) and multiplying \(X,Y\) by the inverse original
coordinate weights. Thus each original coordinate is literally
invariant. The reduced root maps include the original A1 or A2
resolution.

Here are the full point-line correspondences used below. At a reduced
root the point \(aG'\) has identity component. Both the original line
and \(N=\mathcal O(aG'-O)\) have zero degrees on every component. On
the resolved fibre product the first line has zero degree on every
curve contracted by relative minimalization. The blowup Picard
decomposition therefore descends it. Its ratio to \(N\) has vertical
divisor with zero component intersections; the III or IV intersection
matrix has primitive integral kernel the whole-fibre vector. The
actual base equation, with its integer exponent and deck character,
is the resulting Cartier comparison. Its reciprocal is the inverse.

For the good IV-star root the minimal cubic quotient has tangent
weights \((2,1)\) at its three fixed points. On the common root
refinement use rays
\[
 (1,0),(2,1),(1,1),(1,2),(0,1).
 \tag{RJC.19}
\]
For a primitive ray \(v=(v_\eta,v_z)\), the difference between the
root section and the pulled-back resolved section is
\(\min(v_\eta,2v_z)\). Its three coefficients are \(2,1,1\).
Subtracting the \(O\) contribution from the \(aG'\) contribution
gives the complete F46.J.1 Cartier equality. At order \(3l\), pull
the entire diagram back by \(w=z^l\) and normalize its common
blowups. The coefficients are now
\[
 D_P(v)=\min(v_\eta,2l v_z),\qquad
 \beta^*N=\pi^*\mathcal M+D_{aG'}-D_O.
 \tag{RJC.20}
\]
Every toric chart map sends its two monomials to the original
\((\eta,z)\); inversion of the unimodular cone matrix gives its
inverse on the common torus. The strict-section function remains
in the Cartier numerator. This specifies the original point module
as well as every exceptional factor.

For the minimal good III root \(w\), the action is
\[
 (w,X,Y)\longmapsto(iw,-X,iY).
 \tag{RJC.21}
\]
It fixes \(O,T=(0,0)\) and has the order-two orbit
\((i,0),(-i,0)\). Blow up those four points and quotient. The
exceptional images have squares \(-4,-4,-2\); the central image
has square \(-1\), and the fibre is \(4C+D_0+D_1+2D_2\).
Contract \(C\), then the resulting \((-1)\)-curve \(D_2\).
The remaining curves have squares \(-2\) and intersection two.
The invariant substitution in RJC.5 gives the actual III
tangency, with its length-two intersection scheme. The marked
sections avoid both contractions. Consequently
\[
 \beta^*N=\pi^*\alpha^*\mathcal M+E_{aG'}-E_O.
 \tag{RJC.22}
\]
For order \(4k\) the pullback \(w=z^k\) resolves the ideals
\((\eta,z^k)\) at \(O,T\). A complete regular fan has rays
\((j,1)\), \(0\le j\le k\), and \((1,0)\). On successive cones its
literal maps are
\[
 \eta=u^jv^{j+1},\quad z=uv\quad(0\le j<k),\qquad
 \eta=uv^k,\quad z=v
 \tag{RJC.23}
\]
on the last cone. The cone determinants are \(-1\), so their inverse
monomials are integral on the common torus. Pulling the ideal through
these maps gives exactly \(\min(v_\eta,k v_z)\); retaining these
generators in RJC.22 gives its higher-order comparison. The two
remaining fixed-point charts retain their source blowups and carry
no marked point divisor.

For a reduced III root \(h=4k+1\), the original formal parameter
\(\tau=-x/y\) pulls back as \(\tau'=z^k\tau\). The actual
\(aG'\) graph, for \(a=2,6\), is
\(\tau'=z^k A_a(z^h)\), where \(A_a(0)=a\).
The normalized blowup of \((\tau',z^k)\) carries the exact two
Cartier ratios
\[
 \frac{p-A_a(z^h)}p,\qquad 1-vA_a(z^h),\qquad
 \tau'=z^kp,\quad z^k=\tau'v.
 \tag{RJC.24}
\]
The intermediate regular charts are the monomial charts RJC.23.
The source A1 chart is
\[
 X=z\alpha,\quad Y=z\beta,\quad
 \beta^2=z\alpha^3+\alpha+z^{2k}.
 \tag{RJC.25}
\]
At \(k=0\), this includes the original equation with the constant
term \(1\). The inverse is division by \(z\) off the exceptional
curve, and the equation supplies its regular extension. Translation
by \(\pm G'\) uses the actual chord slopes
\((Y\mp z^{k+1})/X\). Resolving their graph gives the mutually
inverse translations of the minimal regular surface. Their
point-module maps carry RJC.24 to the translated point graphs.

The reduced IV root uses the A2 charts of R0E.60-\/-63: in the first
chart \(X=zX_1,Y=zY_1\), the two nonidentity labels are \(Y_1=\pm1\).
The original coordinate substitutions in RJC.4 identify these
labels and the whole exceptional chain. The point graph is the
full ratio \(\tau'/(\tau'-z^{2l+1}L_a(z))\). For the actual
\(3G\) point in RJC.26 its original-coordinate substitution gives
\[
 L_3(z)=-\frac{64+8z^m}{512+96z^m+3z^{2m}},
 \qquad L_3(0)=-1/8.
\]
For the actual doubled point \(6G\) in the same formulae,
\(L_6(0)=-1/4\). The full rational function
\(L_a=-z^{-(2l+1)}X_{aG'}/Y_{aG'}\)
is retained in its own point graph.
Its reciprocal and the inverse chord map give the reverse line
and source arrows. These are the complete F46 and R0E diagrams
for the specific four rows, with their full Cartier functions.

\par

\medskip\noindent\textbf{RJC.2. The ordinary lift, exponents and exact root characters}

\par\nobreak

Retain the rational point-line lift \(U\) of \(+G\) obtained from the
actual source Miller maps. The chosen original forward Hom section is
\[
 B_0=C_a\,\frac{q}{t^3}\,U^{-b},\qquad C_a\in\mathbb C^*.             \tag{RJC.6}
\]
Here \(C_a\) is the coefficient of that specific section in these
specific Miller frames, fixed by evaluation at the same ordinary
reference point; it remains in every map. We verify RJC.6 on the
original regular surface. Put \(T(p)=p-G\). The tangent and chord
formulae give
\[
 \begin{split}
 f_1&=1,\qquad
 f_2=\frac{y-t-x/2}{x-1/4},\\
 f_3&=\frac{(x/4)(y-t-x/2)}
 {(y+t+x/2)/4-(1/4+2t)x},\\
 (x_3,y_3)&=(8t+64t^2,\,3t+96t^2+512t^3),\\
 L_3&=(3x_3^2+t)/(2y_3),\\
 x_6&=L_3^2-2x_3,\quad y_6=L_3(x_3-x_6)-y_3,\\
 f_6&=f_3^2\frac{y-y_3-L_3(x-x_3)}{x-x_6}.
 \end{split}
 \tag{RJC.26}
\]
Their divisors are \(a(G)-(aG)-(a-1)(O)\). Therefore
\[
 B_{a,1}(p)=f_2(p)^a\frac{f_a(Tp)}{f_a(p)},\qquad
 U(p,v)=\bigl(p+G,B_{a,1}(p+G)^{-1}v\bigr)
 \tag{RJC.27}
\]
are the actual negative one-step point-line map and its positive
inverse. The coefficient of the inverse is evaluated at \(p+G\).
All zero and pole factors in RJC.26 belong to the corresponding
point Cartier modules, so these identities are Hom identities
before any specialization.

On the original III components the orders of \(f_2,f_3,f_6\) are
\((0,1),(0,1),(0,3)\). To see these in the original coordinates,
use \(x=p^{-2},y=-p^{-3}\) on the identity component and
\(x=t\alpha,y=t\beta\),
\(\alpha+t\alpha^3=\beta^2-1\), on the other.
The translation \(T\) has directed component formulae
\(\beta'=p-1,p'=\beta-1\).
Substitution into the entire functions RJC.26 gives the orders
just stated and preserves the point factors of R0R.11.
It follows that the two component orders of \(B_{a,1}\) are
\[
 \begin{array}{c|cccc}
 a&1&2&3&6\\ \hline
 \operatorname{ord}_{(C_0,C_1)}B_{a,1}
 &(0,1)&(1,1)&(1,2)&(3,3).
 \end{array}
 \tag{RJC.28}
\]
The \(b\) chronological factors of \(U^{-b}\) thus have order
three on each original component in all four rows. There is no
other finite whole-fibre divisor: the two substitutions in
\(f_a(Tp)/f_a(p)\) have equal generic orders on an irreducible
fibre, and \(f_2\) has zero whole-fibre order at every finite
nonzero \(t\), including the selected node.
The exact integral source translation in F46.J.1,J.6 gives the
Hom class \(F\). Its only remaining whole-fibre coefficient, at
infinity, is consequently \(-2\): \(3F-2F=F\).
Multiplication by \(q/t^3\) leaves exactly the chosen simple
whole-fibre zero at \(t_c\). Its ratio to the original specified
section \(B_0\) has zero divisor on the compact regular source.
Both it and its reciprocal are regular and hence constant.
Reference evaluation defines precisely the nonzero \(C_a\) in
RJC.6. This proves that identity for all four actual point lines.
Its literal rational norm identity is
\[
 \rho=\frac{B_0^{-1}}{U^b}=\frac{t^3}{C_aq}.                         \tag{RJC.7}
\]
For completeness, the new odd-line good-root coefficients can be
read directly before taking central values. At the finite good root
put \(w=z^k\). Then \(f_{2,\rm old}=w^2\phi\),
\(f_{3,\rm old}=w^2\psi\), where
\[
 \phi=\frac{wY-w^2-X/2}{w^2X-1/4},\qquad
 \psi=\frac{X(2wY-2w^2-X)}
 {2wY+2w^2-X-16w^4X}.
 \tag{RJC.29}
\]
The central negative source translation is
\((X,Y)\mapsto(1/X,-Y/X^2)\).
Thus the normalized negative coefficients for \(a=1,3\) are
\(2X,8X\). Their positive inverse coefficients, evaluated at the
inverse source point, are \(X/2,X/8\). The negative twofold
norms are \(4,64\), and the \(a=1\) sixfold norm is \(64\).
These are values of the full functions RJC.29 in the full
Cartier comparison RJC.22-\/-23.

At the good infinity root put \(v=z^l\). The full coefficient is
\[
 f_{2,\rm old}=v^{-1}
 \frac{Y-1-vX/2}{X-v^2/4}.
 \tag{RJC.30}
\]
On the central cubic the positive source map is
\[
 X_+=2(1-Y)/X^2,\qquad
 Y_+=(Y-1)(X-X_+)/X-Y.
 \tag{RJC.31}
\]
Its \(a=1\) inverse-source line coefficient is \(X/2\);
its three chronological factors multiply to \(-1/2\) and
its six factors to \(1/4\). The opposite sixfold norm is \(4\).
Together with RJC.20 these formulae carry the odd point line
through its full original infinity comparison.

Put
\[
 J_0=\frac{mh}{b},\qquad
 d_0=\frac{3\varepsilon-mn_0}{b}
     =\frac{2\varepsilon+hn_\infty}{b}.                              \tag{RJC.8}
\]
Both are integers on every domain RJC.2. The equality follows by
subtracting RJC.3. Substituting that same actual equation gives the
explicit integral expressions
\[
\begin{array}{c|c|c}
(a,b)&J_0&d_0\\ \hline
(1,6)&2kl&2kn_\infty+ln_0\\
(2,3)&l(4k+1)&(4k+1)n_\infty+2ln_0\\
(3,2)&2k(3l+2)&6kn_\infty+(3l+2)n_0\\
(6,1)&(3l+2)(4k+1)&3(4k+1)n_\infty+2(3l+2)n_0.
\end{array}                                                        \tag{RJC.9}
\]

The actual extending normalized root lift \(L_i\) in K67.SE.3 satisfies
\(L_i^b=B_0^{-1}\). The line correspondences and that norm identity
retain the following source comparison:
\[
\begin{array}{c|c|c}
i&e_i&U=z^{e_i}u_i(z)L_i,\quad u_i^b\\ \hline
0&-3h/b&C_aq\\
\infty&2m/b&C_a h_\infty,\quad h_\infty=1+4z^m/27.
\end{array}                                                        \tag{RJC.10}
\]
The constants of the two analytic unit roots are fixed by the actual
source/Cartier correspondence. Here is the extension and normalization
that produces them. The root chord lift of \(+G'\) is a meromorphic
Hom section. Its horizontal divisor is the prescribed difference of
point modules. After the exact point corrections RJC.20-\/-25 its
remaining vertical divisor has zero component degrees. The primitive
kernel of the resolved intersection matrix makes it an integer
whole fibre. Divide by that actual power of \(z\).
The resulting section and its inverse have zero orders at every
codimension-one point. Normal local rings express each as the
intersection of their height-one valuation rings; hence both are
regular, including at intersections and section points. They give
the extending positive lift with its extending inverse.

Its deck commutator is a nonzero base unit. Choose its logarithm
with the prescribed value at the reference. For its norm-zero
nonconstant part \(g\), the finite averaging
\(-m_i^{-1}\sum_{j=0}^{m_i-1}jD_i^jg\) supplies the correcting
logarithm. An invariant analytic \(b\)-th root of the remaining
norm unit then gives \(L_i^b=B_0^{-1}\). Both corrections are
actual scalar maps on the root line, with their reciprocal maps.
The finite value of the commutator is retained as \(\chi_i\).
Comparing the norm to RJC.7 fixes exactly the integer powers
and analytic units in RJC.10. In particular these units are
retained, and their reciprocals give the inverse comparisons.
Their exponent identity is \(\rho=z^{-be_i}u_i^{-b}\).

The deck characters are the actual fibre-weight ratios:
\[
\begin{array}{c|c|c}
(a,b)&\chi_\infty&\chi_0\\ \hline
(1,6)&\zeta_3&-1\\
(2,3)&\zeta_3^2&1\\
(3,2)&1&-1\\
(6,1)&1&1.
\end{array}                                                        \tag{RJC.11}
\]
They give, with the complete cap generator \(A_{i,r}\),
\[
 L_i^\varepsilon A_{i,r}L_i^{-\varepsilon}
 =\chi_i^{-\varepsilon}A_{i,r}
 =\zeta_b^{Mn_i/m_i}A_{i,r}
 =A_{i,\zeta_b r}.                                                  \tag{RJC.12}
\]
The middle equality is exactly RJC.3 modulo two and three. For example,
in the first row \(kn_\infty\equiv\varepsilon\pmod3\) and
\(ln_0\equiv\varepsilon\pmod2\). The first congruence gives
\(\zeta_6^{4kn_\infty}=\zeta_3^{-\varepsilon}\), and the second gives
\(\zeta_6^{3ln_0}=(-1)^{-\varepsilon}\). The other rows have the same
displayed reductions with the trivial partner character retained.
Every prescribed common finite line character cancels in this
conjugation.

\par

\medskip\noindent\textbf{RJC.3. The full ordinary endpoint and its cyclic norm}

\par\nobreak

Write the actual one-gamma forward closing as
\[
 h_0(t)=c_0 q^\varepsilon t^{-mn_0},\qquad c_0\in\mathbb C^*,          \tag{RJC.13}
\]
with its original annular-reference constant \(c_0\).
Choose the prescribed scalar \(\gamma_0\) with
\(\gamma_0^b=c_0C_a^{-\varepsilon}\).
The rational root required by the full K67 comparison is literally
\[
 g_0=\gamma_0t^{d_0},\qquad
 g_0^b=h_0\rho^\varepsilon.                                         \tag{RJC.14}
\]
On \(\lambda=r^M\) the ordinary positive endpoint is therefore
\[
 E_r=\gamma_0 r^{-\varepsilon J_0}t^{d_0}U^\varepsilon.               \tag{RJC.15}
\]
It covers \(+\varepsilon G\). Its inverse evaluates \(U^{-\varepsilon}\)
at the actual inverse source point and uses the reciprocal displayed
base and scalar factors. At the finite root its total \(z\)-power is
\[
 hd_0-3h\varepsilon/b=-J_0n_0,
\]
and at infinity it is
\[
 -md_0+2m\varepsilon/b=-J_0n_\infty.                                 \tag{RJC.16}
\]
The full forward ordinary-to-cap cochain for the path winding \(J_0\)
contributes \(z^{J_0n_i}\). Thus the remaining cap arrow is
\(\gamma_0 r^{-\varepsilon J_0}u_i^\varepsilon L_i^\varepsilon\),
including the full invariant unit and finite character in RJC.10-\/-12.
The full annular terms and all original Cartier generators in
this square are given in RJC.4 below.

The \(b\) chronological source points of the parameter are
\(r,\zeta_b r,\ldots,\zeta_b^{b-1}r\). Their scalar product is
\[
 r^{-\varepsilon bJ_0}
 \exp\!\left(-\pi i\varepsilon J_0(b-1)\right).
\]
For all four rows \(J_0(b-1)\) is even: the even \(b\) rows have even
\(J_0\), and the odd \(b\) rows have even \(b-1\). Keeping the
reference constants, RJC.6,9,13 therefore give the exact ordinary
composite
\[
 E_{\zeta_b^{b-1}r}\cdots E_{\zeta_b r}E_r
       =h_0(t)B_\lambda^{-\varepsilon}.                             \tag{RJC.17}
\]
The K67 bounded origin has the additional constant
\[
 c=\exp\!\left(\frac{2\pi i\varepsilon M(b-1)}{2b^2}\right).
\]
Its exact family comparison is the physical scalar
\[
 F_r=r^d,\qquad d=\frac{\varepsilon J_0(b-1)}2\in\mathbb Z,
 \qquad F_{\zeta_b r}E_r=(cE_r)F_r.                                 \tag{RJC.18}
\]
The inverse is \(r^{-d}\). This keeps the finite phase even when \(c\)
is a nontrivial root of unity; its \(b\)-fold product is already the
same exact closing in RJC.17.

The next sections prove that this is a family comparison through
the complete roots, references and fans and that its powers
exhaust all holomorphic endpoints on this domain.

\par

\medskip\noindent\textbf{RJC.4. Both complete annular references and their scalar variation}

\par\nobreak

At site \(i\in\{0,\infty\}\), write \(q_0=t\),
\(q_\infty=1/t\), \(m_0=h,m_\infty=m\).
The actual root division section \(R_i\) satisfies
\(m_iR_i=bG'\). At a reduced root multiplication by \(m_i\)
is invertible on both the additive identity group and the finite
component group. At a good root use its analytic inverse branch
at \(O\). These give the unique marked analytic division section,
and the deck preserves it. Its negative translation is the
source part of the cap lift.

For \(P_i=aG'\), the chord-over-vertical function
\[
 g_{P_i,R_i}(p)=
 \frac{\ell_{P_i,R_i}(p)}{x(p)-x(P_i+R_i)}
 \tag{RJC.32}
\]
has divisor \((P_i)+(R_i)-(P_i+R_i)-(O)\).
It is precisely the point-line Hom coefficient over \(p\mapsto
p-R_i\). Its whole-fibre power is removed by the actual graph
generators RJC.20-\/-25; the full strict-section factors remain.
Finite averaging as in RJC.2 supplies the deck-equivariant lift.
For this division lift the remaining deck character is one.
At a good root \(R_i(0)=O\), and its central line map is a
constant scalar; its commutator with the deck is one there.
At a reduced root of order greater than one, the explicit
graphs RJC.24 and the reduced-IV graph show \(aG'(0)=O\).
The central line is trivial with its actual point/O
identification, so the central translation lift is again
a constant in that identification. At order one the deck
is identity. In each case the constant character left by
finite averaging equals its central value one.
The ratio of the \(m_i\)-fold norm to \(B_0\) is therefore an
invariant nonzero analytic base function. Taking the root fixed
by the prescribed reference gives the full commuting lift
\(\widetilde F_i^{\,m_i}=B_0\). Thus the
actual marked cap generator, including its prescribed character,
is
\[
 A_{i,\lambda}=R_{\lambda,i}^{\,n_i}
 \zeta_{m_i}^{d_i}\widetilde F_i^{\,n_i}D_i,\qquad
 R_{\lambda,i}^{\,m_i}=\lambda,\qquad
 A_{i,\lambda}^{\,m_i}=B_\lambda^{\,n_i}.
 \tag{RJC.33}
\]
The finite character \(d_i\), division-section branch and scalar
root are those of the original cap. The source and line maps
commute with \(B_\lambda\). In constant-period coordinates they
are translations in the commutative extension
\((\mathbb C\times\mathbb C^*)/\Lambda\); the Cartier extensions
above carry this identity through the whole root.

On the reference sectors use
\[
 s_{i\alpha}=\log_\alpha q_i/(2\pi i),\quad
 z_{i\alpha}=e^{2\pi i s_{i\alpha}/m_i},\quad
 c_{\lambda,i}=
 n_i\log R_{\lambda,i}/(2\pi i)+d_i/m_i,\quad
 f_{i\alpha}=\operatorname{Log}_\alpha\widetilde F_i.
 \tag{RJC.34}
\]
With the full original-coordinate and Cartier map \(I_{z_{i\alpha}}\)
on the annulus, the gluing arrow is
\[
 \Phi_{i\alpha}=
 \operatorname{Exp}
 \bigl(n_i s_{i\alpha}f_{i\alpha}
       +s_{i\alpha}c_{\lambda,i}\delta+h_{i\alpha}\bigr)
 I_{z_{i\alpha}}.
 \tag{RJC.35}
\]
It includes the negative source displacement
\(-n_i s_{i\alpha}R_i\), every point-line factor, and the full
physical scalar. On a labelled overlap
\(s_{i\beta}=s_{i\alpha}+k\),
\(f_{i\beta}=D_{i*}^kf_{i\alpha}+\Pi_{i\alpha\beta}\).
The complete cochain equation is
\[
 h_{i\beta}-D_{i*}^k h_{i\alpha}
 =-n_i s_{i\alpha}\Pi_{i\alpha\beta}
  +j_{i\alpha\beta}\operatorname{Log}_\beta B_\lambda
  +\Lambda_{i\alpha\beta}.
 \tag{RJC.36}
\]
The source periods \(\Pi,\Lambda\) and integral suspension
columns \(j\) remain those of the specified reference.
An explicit solution in its ordinary Lie frame starts with its
fixed smooth extension \(h^0\), its residual cocycle \(E\) and
a partition \(\chi_\beta\):
\[
 k_\alpha=-\sum_\beta\chi_\beta E_{\alpha\beta},\qquad
 v(q)=\frac1\pi\int
 \frac{\bar\partial k(w)}{q-w}\,dA(w),\qquad
 h_\alpha=h_\alpha^0+k_\alpha-v.
 \tag{RJC.37}
\]
The residual is supported on the prescribed compact subannulus.
The common holomorphic constant of \(v\) is fixed by the original
reference value; equivalently subtract \(v(q_{\rm ref})\) and
apply the corresponding common constant to the rigidification.
This gives exactly that value and leaves every integral period
in RJC.36 unchanged. Subtracting the target-sheet logarithm
\(kn_i f_{i\beta}+kc_{\lambda,i}\delta\) proves
\[
 \Phi_{i\beta}=
 A_{i,\lambda}^{\,k}B_\lambda^{\,j_{i\alpha\beta}}\Phi_{i\alpha},
 \qquad
 \Phi_{i\alpha}B_\lambda=B_\lambda\Phi_{i\alpha}.
 \tag{RJC.38}
\]
The inverse is the inverse root map followed by the negative
of the displayed Lie translation at the inverse source point,
with the reciprocal point and scalar coefficients.

For scalar variation put
\(\lambda_\theta=e^{2\pi i\theta}\lambda_*\), and keep the
actual integer representatives \(J_{i\alpha}\) with
\(J_{i\beta}-J_{i\alpha}=j_{i\alpha\beta}\).
The full reference restoration is
\[
 \Phi_{\theta,i\alpha}
 =e^{2\pi i n_i\theta s_{i\alpha}/m_i}
 B_\theta^{\,J_{i\alpha}}
 \bigl(B_*^{-J_{i\alpha}}\Phi_{*,i\alpha}\bigr).
 \tag{RJC.39}
\]
On an integer winding \(J\), the last two factors have ratio
\(e^{2\pi iJ J_{i\alpha}}=1\), and the first factor gives
\[
 \Phi_{\theta+J,i\alpha}
 =z_{i\alpha}^{\,n_iJ}\Phi_{\theta,i\alpha}.
 \tag{RJC.40}
\]
This establishes the root factors used in RJC.16 with every
source period and line reference still present. The remaining
cap map is the complete \(u_i^\varepsilon L_i^\varepsilon\)
map of RJC.10, multiplied by the specified scalar.
RJC.12 supplies its finite generator square.

The integral form of that same square is the K67 cochain
\[
 \frac{P_i\delta-\varepsilon(A_i-I)e_\gamma}{b}
 =\bigl(0,-\varepsilon k_{i,2},\varepsilon k_{i,1},
                (P_i+\varepsilon z_i)/b\bigr)^t,\qquad
 P_i=Mn_i/m_i.
 \tag{RJC.41}
\]
Here \(z_i=2\) at infinity and \(-3\) at zero are representatives
of the actual characters \(\chi_i=\zeta_b^{z_i}\).
RJC.8-\/-12 make the last coordinate integral, and the source
coordinates are the original integral Abel-period columns.
This exhibits the entire ordinary/root lattice arrow.

\par

\medskip\noindent\textbf{RJC.5. The intrinsic base and full affine endpoint map}

\par\nobreak

The integral chain maps F46.J.7-\/-12 give the actual sphere
topology under RJC.3: their degree-two map is an integral
isomorphism and both degree-one and degree-three determinants
are \(-\varepsilon\). The cap fundamental-group presentation
reduces to the cyclic group of order \(|\varepsilon|=1\).
The common blowups and both ruling maps of F35 carry these
integral maps through each prescribed fan. In particular every
object here has \(H^2=H^4=0\).

Let \(f\) be a meromorphic function on one of these compact
objects. A horizontal component of its polar divisor restricts
to a nonempty effective divisor on a generic compact two-torus
and has positive pairing with a flat Kähler form there.
Its cohomology class is the restriction of a class in
\(H^2(X)=0\), a contradiction. All poles are vertical, and
restriction to each regular compact torus is constant.
The original transverse maps at special components are
\(u\mapsto t_0+u^e\). Their invariant Laurent expansions
contain precisely multiples of \(e\); hence the constants
descend meromorphically across each special value. We obtain
the actual field isomorphism
\[
 \pi^*:\mathbb C(t)\xrightarrow{\;\cong\;}\mathbb C(X).
 \tag{RJC.42}
\]
Its inverse is restriction to a regular fibre value followed
by the just-described Laurent descent. A holomorphic endpoint
and its inverse consequently induce inverse Möbius maps of
the original base.

The complete generic point-line period module is, for all four
rows,
\[
 \gamma=(-bg,\ell),\quad e_1=(1,0),\quad
 e_2=(\tau,ag),\quad\delta=(0,1).
 \tag{RJC.43}
\]
The full theta/Miller reference constant occurs in \(\ell\);
both point Cartier modules are included in these coordinates.
For a source loop \(\alpha\) its monodromy is
\[
 M_\alpha=
 \begin{pmatrix}
 1&0&0\\
 b v_\alpha&A_\alpha&0\\
 c_\alpha&a v_\alpha^t J A_\alpha&1
 \end{pmatrix},\quad
 J=\begin{pmatrix}0&-1\\1&0\end{pmatrix},
 \tag{RJC.44}
\]
with the \(1,2,1\) block order. Direct multiplication gives
\(v_{\alpha\beta}=v_\alpha+A_\alpha v_\beta\) and
\[
 c_{\alpha\beta}=c_\alpha+c_\beta+
 ab\,v_\alpha^tJ A_\alpha v_\beta.
 \tag{RJC.45}
\]
Thus no central cochain has been removed.

The Abel cocycle \(v_\alpha\) is nonzero over \(\mathbb Q\).
A rational coboundary, after clearing its denominator and the
finite component orders, would give a global holomorphic Lie
logarithm of a nonzero multiple of \(G\). On the minimal good
III chart the logarithm of \(2G'\) divided by its actual
factor \(w\) extends; the order-four elliptic action has no
invariant period. On the minimal good infinity chart
\(Y^2=X^3+wX+1\), the logarithm of \(3G'\) has leading term
\(-w^2/8\); division by \(w^2\) gives the original extending
Lie coefficient. Its order-three elliptic action has no
invariant period. At the selected node the invariant annular
period extends through \(xy=q\).
The original Lie transition is
\[
 u_\infty(q_\infty)
 =q_\infty^{-1}u_0(q_\infty^{-1}).
 \tag{RJC.46}
\]
An entire \(u_0\) satisfying this transition tends to zero
at infinity and is zero. But \(3G\) has original infinity
value \((64,512)\) and additive parameter \(-1/8\);
each of its nonzero multiples has nonzero additive parameter.
This contradiction proves the asserted Abel obstruction.
The higher-root diagrams RJC.20-\/-25 are the pullbacks of
these actual minimal diagrams, so the argument applies to
every cap order in RJC.2.

The vector and dual equations of RJC.44 now give
\[
 \Lambda^{\pi_1}=\mathbb Z\delta,\qquad
 (\Lambda^*)^{\pi_1}=\mathbb Z\psi,\quad
 \psi(\gamma)=1,\quad\psi(e_1)=\psi(e_2)=\psi(\delta)=0.
 \tag{RJC.47}
\]
A nonzero \(\gamma\) coefficient in an invariant vector makes
the Abel cocycle a coboundary. With that coefficient zero,
the order-three infinity matrix kills the remaining invariant
elliptic vector. The dual equation has the same nonzero Abel
obstruction. These are primitive integral subgroups, so the
extra deck of the kernel cover is the actual primitive
\(\gamma\). The kernel cover is the nonzero point line
with its actual Cartier modules.

An endpoint preserves these two subgroups with their signs.
The derivative of its universal-cover lift is periodic for
the full compact torus lattice and therefore constant.
Complex linearity and \(\delta\mapsto e\delta\) give the
full affine map and its inverse
\[
 \begin{split}
 (z,w)&\longmapsto(Az+z_0,Cz+e w+w_0),\\
 (z',w')&\longmapsto
 \bigl((z'-z_0)/A,\,
 e(w'-w_0-C(z'-z_0)/A)\bigr),
 \quad A\ne0,\quad e=\pm1.
 \end{split}
 \tag{RJC.48}
\]
On the kernel lattice its complete integral matrix and inverse
are
\[
 K=\begin{pmatrix}p&q_1&0\\r&s&0\\u&v&e\end{pmatrix},
 \quad
 K^{-1}=
 \begin{pmatrix}
 s&-q_1&0\\-r&p&0\\e(vr-us)&e(uq_1-vp)&e
 \end{pmatrix},\qquad ps-q_1r=1.
 \tag{RJC.49}
\]
The positive determinant follows from the oriented area of
the nonzero complex-linear first coordinate. The first
coordinates of RJC.48 are actual inverse elliptic maps.
For this source
\[
 j_0(t)=\frac{6912}{27t+4}.
 \tag{RJC.50}
\]
The equality \(j_0(h(t))=j_0(t)\) forces \(h(t)=t\).
The endpoint therefore fixes the original base everywhere,
including its two original cap centres and selected height.

The full linear centralizer of RJC.44 is \(eI+c\delta\psi\).
For the elliptic block, the infinity matrix generates
\(\mathbb Q(\zeta_3)\); the selected unipotent has a nonzero
nilpotent part. A nonscalar element of that quadratic field
generates the field, whose nonzero elements are invertible.
Its simultaneous centralizer with the unipotent is thus
\(\mathbb QI\). The two off-diagonal block equations use
the nonzero Abel class to identify the adjoining scalars
with this one and to set their elliptic vectors to zero.
The central cochains remain in the calculation and leave
exactly \(c\delta\psi\). Integral unimodularity gives
\(e=\pm1\).
The two complex-independent period vectors \(e_1,\delta\)
force the derivative to be exactly \(eI\). The \(\gamma\)
equation then says that the two suspension logarithms
differ by an integer. Thus an endpoint has exactly equal
original scalar values at its two ends.

\par

\medskip\noindent\textbf{RJC.6. Extending the source displacement and the full section lattice}

\par\nobreak

Consider first a positive endpoint. Its ordinary lift has the
form \(H(t)U^j\) after extending its source displacement.
We give the extension rather than assume it. Pull the cap
map and its inverse to the normalized fibre product
\(z_i^{m_i}=q_i\). On the actual point-line cover, a
nonzero vector over the \(O\)-section kills the root meridian;
the remaining vanishing cycles lie in \(\ker\psi\).
These are the Cartier charts RJC.20-\/-25, so the lift is
single-valued on the full root source.
Its commutator with the cap generator is an integer power
\(B^s\). Taking the \(m_i\)-fold relation gives
\(B^{m_i s}=1\). The original strict contraction has no
nonzero finite-order deck, hence \(s=0\).
The source displacement is therefore an equivariant
holomorphic root section. The original-coordinate maps
RJC.4-\/-5 descend it, and the proper minimal regular
source resolves the indicated graph.

At the selected fibre a pole of order \(-K\) in the
physical coefficient is cancelled by the actual \(B^K\)
reindex. Its remaining source displacement is shifted by
\(-bKG\), and its scalar coefficient is a nonzero
holomorphic torus unit. In the two selected vertex
charts these maps are the unit maps displayed in
RJC.65 below, with reciprocal inverses.
They extend through each vertex and open curve axis.
The \(B\)-labels of two ordinary local lifts differ by
a locally constant integer. After the selected unit
has been joined, their joining torsor is over the
two-point root sphere with presentation
\[
 \langle x,y\mid x^m=y^h=xy=1\rangle.
 \tag{RJC.51}
\]
This group is cyclic of order \(\gcd(m,h)=1\).
Its integer joining torsor is trivial, with the
constant join fixed at the original reference.
The source displacement extends to an actual global
section of the original elliptic surface.

The section group here is exactly \(\mathbb ZG\).
To verify the integral index, use the birational chart
\[
 u=\frac{y-t-x/2}{x},\quad
 v=\frac{y+t+x/2}{x},\quad
 x=uv+\frac14,\quad
 t=\frac{x(v-u-1)}2,\quad
 y=\frac{x(u+v)}2.
 \tag{RJC.52}
\]
The displayed reciprocal functions are inverse on their
common open. The source is a rational elliptic surface
with Euler number \(8+3+1=12\).
Its ten classes
\[
 (O,F,C_0,\ldots,C_5,A_1,G)
 \tag{RJC.53}
\]
have \(O^2=G^2=-1\), \(OF=GF=1\),
the negative \(E_6\) Cartan block with edges
\((0,1),(1,2),(2,3),(3,4),(2,5)\), and \(A_1^2=-2\).
The additional nonzero intersections are
\(G C_0=G A_1=1\); all unspecified ones are zero.
Their determinant is \(-1\).
They span the entire rank-ten unimodular source lattice.
Quotienting by its first nine classes is free cyclic,
with generator the image of \(G\). This is the full
section quotient, with no remaining torsion or index.
Its height is
\(2-4/3-1/2=1/6\).
Thus the extended displacement above is \(jG\) for
an integer \(j\), and its full ordinary point-line
map is precisely \(H(t)U^j\), up to its physical
scalar already included in \(H\).

\par

\medskip\noindent\textbf{RJC.7. All positive endpoints and the physical automorphism group}

\par\nobreak

For one object, the exact zero and infinity comparisons
RJC.10 and the selected unit reindex give
\[
 \operatorname{ord}_0H=3j/b,\quad
 \operatorname{ord}_\infty H=-2j/b,\quad
 \operatorname{ord}_{t_c}H=-K.
 \tag{RJC.54}
\]
There are no other zeros or poles: the map and inverse
are nonzero Hom arrows on every other fibre.
The degree of this principal divisor gives \(j=bK\),
and \(H=c(t^3/q)^K\).
Using the exact \(C_a\) in RJC.6, the quotient map is
the physical scalar
\[
 H U^{bK}=cC_a^K B_0^{-K}
        =cC_a^K\lambda^K B_\lambda^{-K}
        \quad\text{on the covering line.}
 \tag{RJC.55}
\]
All physical scalars extend by the displayed root,
Cartier and selected unit maps. A scalar acting as an
identity deck would translate the source by a multiple
of \(bG\); the nonzero Abel class shows that \(G\) has
infinite order, so that multiple is zero and the scalar
is one. The physical action is faithful.

A negative automorphism of this one object has
sixteen transverse fixed points on each regular
two-torus, because its affine linear part is \(-I\).
Its compact analytic equalizer contains a horizontal
curve of positive base degree. The degree equals
pairing that curve with \(\pi^*H^2(\mathbb P^1)\),
whose image in \(H^2(X)\) is zero. The contradiction
excludes every negative automorphism. We have the
actual group isomorphism, with its explicit inverse
from RJC.55,
\[
 \mathbb C^*_{\rm physical}
 \xrightarrow{\;\cong\;}\operatorname{Aut}_{\rm hol}(X).
 \tag{RJC.56}
\]

For a positive endpoint of scalar winding \(J\), RJC.40
contributes both root factors. Its full divisor equations
are
\[
 \begin{split}
 \operatorname{ord}_0H&=\frac{3j}{b}-\frac{n_0J}{h},\\
 \operatorname{ord}_\infty H&=-\frac{2j}{b}
                                  -\frac{n_\infty J}{m},\\
 \operatorname{ord}_{t_c}H&=-K.
 \end{split}
 \tag{RJC.57}
\]
These are integer orders of the single-valued ordinary
base function; all corresponding local units were retained
in RJC.10,35-\/-40. Their degree and RJC.3 give
\[
 \frac jb-K-\frac{\varepsilon J}{mh}=0,\qquad
 j=bK+\frac{b\varepsilon J}{mh}.
 \tag{RJC.58}
\]
It follows that \(M\mid bJ\), or \(J\in J_0\mathbb Z\).
The endpoint RJC.15 attains \(J_0\). Its chronological
powers and inverse-source inverses attain every integer
multiple. At a fixed winding the ratio of two endpoints
is an automorphism RJC.56, hence exactly a physical
scalar with its full extending map.

\par

\medskip\noindent\textbf{RJC.8. Single-valued negative endpoints}

\par\nobreak

The selected inversion
\[
 (Z,V,q)\longmapsto(Z^{-1},V^{-1},q)
 \tag{RJC.59}
\]
reverses both full periods. Negating a lattice-basis
triangle gives its opposite triangle, translated by
\(-p-u-v\). The actual vertex exchange in R0E.90
preserves \(q\), the factors \(\kappa g,\alpha\eta\),
and their reciprocal factors. It therefore extends
on every selected fan, with itself as inverse.

A negative endpoint has single-valued nonzero-line
lifts on the ordinary and root pieces: the point/O
vector and root-cover argument of RJC.6 applies to
its inverse as well. Composing its selected arrow
with RJC.59 gives the extending positive selected
unit arrow, so its continuation has no additional
selected integer monodromy. The constant integral
joins use RJC.51. Both lifts reverse the full
suspension,
\(\widetilde F B=B^{-1}\widetilde F\).

On the original infinity annulus the complete fixed
reference is \(A_\infty B^\nu\), with the actual
retained integer \(\nu\). One root circle consequently
has reference \(B^N\), where
\[
 N=n_\infty+m\nu.
 \tag{RJC.60}
\]
The two scalar endpoints have the same original
scalar and the same reference operators by RJC.39-\/-40.
Choose the one constant deck making their two lifts
equal on a cut annulus. Continuation around the
root circle gives
\[
 B^N\widetilde F B^{-N}=B^{2N}\widetilde F.
 \tag{RJC.61}
\]
Single-valuedness implies \(B^{2N}=1\), and strict
contraction gives \(N=0\).
But \(\gcd(m,n_\infty)=1\) and every allowed \(m\)
is at least two, so \(n_\infty+m\nu\ne0\) for
every integer \(\nu\). This excludes all negative
endpoints, including the quadratic reduced infinity
root. The complete holomorphic winding group on
each stated joint domain is therefore exactly
\(J_0\mathbb Z\).

\par

\medskip\noindent\textbf{RJC.9. The entire parameter-family comparison on every fan}

\par\nobreak

We now carry RJC.18 through the actual selected charts.
Let \(v=\log r/(2\pi i)\), let \(\Gamma\) be the bounded
Tate logarithm of \(+bG\), and retain the full ordinary
closing logarithm
\(H_0=\varepsilon s+U_0(q)\).
Here \(s=\log q/(2\pi i)\); the zero-cap factor
\(-mn_0\log t/(2\pi i)\) and the fixed constant
\(\log c_0/(2\pi i)\) are included in \(U_0\), which is
holomorphic near the selected \(t_c\ne0\).
The full outward suspension period there is
\(\gamma=(\Gamma,C-s-Mv)\).
The bounded endpoint displacement is
\[
 \tau(v)=\frac{H_0\delta+\varepsilon\gamma}{b}
 =\left(\frac{\varepsilon\Gamma}{b},
 \frac{U_0+\varepsilon C-\varepsilon Mv}{b}\right).
 \tag{RJC.62}
\]
Every entry is holomorphic on the whole selected disc.
Put \(\alpha=\varepsilon M(b-1)/(2b^2)\) and
\[
 \sigma(v)=
 \left(\varepsilon\Gamma v,\,
 v(U_0+\varepsilon C)
       -\frac{\varepsilon M}{2}(v^2-v)\right).
 \tag{RJC.63}
\]
Direct subtraction gives
\(\sigma(v+1/b)-\sigma(v)=\tau(v)+\alpha\delta\).
Thus \(\xi=z-\sigma(v)\) is the closed native coordinate
for the forward endpoint \(cE_r\).
For every cut interval,
\[
 C_s(v)=\exp(\sigma(v)-\sigma(v+s)),\quad
 P_j(v)=\exp(\sigma(v)-\sigma(v+j/b)),
 \tag{RJC.64}
\]
and \(C_s(v+j/b)P_j(v)=C_{j/b+s}(v)\) literally.
Their inverse maps have the negative exponents.
At \(j=b\) the product is
\(\exp(-H_0\delta-\varepsilon\gamma)\), the full
one-gamma closing. This equality holds on parameter
neighborhoods, so the first and last collars agree
with all derivatives.

For a lattice basis \(P\in GL_2(\mathbb Z)\), conjugate
the full period matrix by \(P\); all factors
\(\kappa,\alpha_0,g,\eta\) stay in that matrix.
The log characters are \(P^{-1}(z_1,z_2)^t\).
The physical scalar \(F_r=r^d\) of RJC.18 acts on
their multiplicative coordinates by
\[
 A=r^{d(P^{-1})_{12}},\qquad
 D=r^{d(P^{-1})_{22}}.
\]
For these factors, or for the full two factors
of any \(C_s,P_j\), the actual two vertex maps are
\[
 \begin{split}
 (a,b_1,c_1)&\longmapsto
 ((AD)^{-1}a,A b_1,D c_1),\\
 (u,v_1,w_1)&\longmapsto
 (ADu,D^{-1}v_1,A^{-1}w_1).
 \end{split}
 \tag{RJC.65}
\]
Their inverses use \(A^{-1},D^{-1}\). They preserve
\(q=ab_1c_1\) and intertwine every full transition,
because each transition is the specified monomial
with its unchanged source coefficient.
For example the vertical fan \(T_j\) retains
\[
 \kappa_j=\kappa g^j,\quad
 \alpha_j=\alpha_0\eta^j\kappa^{-j}g^{-j^2},\quad
 g_j=g,\quad\eta_j=\eta g^{-j},
 \tag{RJC.66}
\]
and its three full transition rows are
\[
 \begin{array}{c|ccc}
 &u&v_1&w_1\\ \hline
 C_0&\kappa_j\alpha_j b_1c_1&
       \alpha_j^{-1}ab_1&\kappa_j^{-1}b_1^{-1}\\
 C_1&g_j\eta_j b_1c_1&
       \eta_j^{-1}c_1^{-1}&g_j^{-1}ac_1\\
 C_2&a^{-1}&ab_1&ac_1.
 \end{array}
 \tag{RJC.67}
\]
Substitution of RJC.65 gives the same character
factor on each side, with every source unit present.

The corresponding native Cox map is
\[
 (a_0,a_1,c_0,c_1)\longmapsto
 (a_0,Aa_1,(AD)^{-1}c_0,Dc_1).
 \tag{RJC.68}
\]
It preserves the Cox relation and descends to both
native charts. Its reciprocal map proves the
inverse on the exceptional curve and the normal
bundle. Normal rescaling retains the complete
\(q\)-jets: \(q\) has order two in the normal
coordinates of an open curve axis and order three
at a vertex, so the remaining source dependence
is precisely the higher normal jet carried by
that rescaling.
RJC.64 is an equality of both full characters,
and hence of RJC.65 and RJC.68 on every sheet.
The ordinary, cap and Cartier maps are
RJC.15-\/-18 and RJC.35-\/-41.
They prove that \(F_r\), with inverse \(r^{-d}\),
is an isomorphism of the complete parameter
families, including every endpoint and intermediate
seam.

The action \(r\mapsto\zeta_b r\) on the parameter
punctured disc is free. Its quotient is \(w=r^b\),
with \(\lambda=w^{J_0}\).
The complete \(r\)-family is the literal degree-\(b\)
pullback of this shorter family.
This is a covering of the completed orbit pieces:
the physical line, cap identification and stable
tangent bundle pull back by the displayed maps.
Integration of their integral characteristic
classes multiplies by \(b\). This supplies the
integer characteristic comparison used below,
before reduction modulo \(28\).

\par

\medskip\noindent\textbf{RJC.10. The original page, its line and the primitive pair}

\par\nobreak

The actual canonical page over an ordinary binding
value \(\tau\) has source-base curve
\[
 \eta^h-\tau\xi^m=1,\qquad
 t=\tau+\xi^{-m}=\eta^h\xi^{-m}.
 \tag{RJC.69}
\]
The reciprocal functions on its ordinary open are
\(\xi^m=(t-\tau)^{-1}\) and
\(\eta^h=t/(t-\tau)\), with all chosen roots
given by the original cap labels.
The degree is \(M=mh\); the three branch indices
at \(0,\infty,\tau\) are \(h,m,M\).
Riemann-\/-Hurwitz gives genus
\((m-1)(h-1)/2\).
The original source page is the pullback by
these actual maps, followed by the minimal root
models RJC.4-\/-5 and their stated resolutions.

The source translation comparison of the page
uses the rigidified relative Picard evaluation
of O4-\/-O5 in the proved order-one-page construction.
It is worth specifying its map on this source.
On each translation-parameter chart, form the
ratio of the original pulled-back point line
and the translated degree-zero root point line.
On every smooth fibre it is trivial; direct
image and evaluation supply its scalar base line.
The evaluation divisor has zero intersection
with every component by RJC.20-\/-25.
The connected Kodaira intersection matrix
therefore makes it an integral sum of whole
fibres. Division by their actual base equations
gives an invertible evaluation map, with
reciprocal inverse, at singular as well as
smooth fibres.
Rigidifying at \(O\) and pulling back by the
smooth translating section gives the exact
page-line comparison.
Every root character and full delta transition
remains as base-line holonomy.
The only resulting change of the line class
is \(c_1(N)\mapsto c_1(N)+zF\), where
\[
 c_1(N)F=F^2=0.
 \tag{RJC.70}
\]
Thus its square is unchanged. On the binding
collar this is the specified comparison; the
base line extends over the ordinary disc and
has zero boundary square ambiguity by the
full K5 relative gauge formula. The stable
tangent comparison is the same source map.

Let \(e_\infty,e_0\) be the Euler numbers of
the actual minimal root fibres:
\((0,0),(0,3),(4,0),(4,3)\) in the four
rows of RJC.2.
The single original I1 away from the cap
centres has \(M\) preimages counted with
degree, infinity has \(h\) root fibres,
and zero has \(m\) root fibres.
The page height and holomorphic Euler number
are therefore
\[
 H=\frac{Ma^2}{6},\qquad
 \chi=\frac{M+h e_\infty+m e_0}{12}.
 \tag{RJC.71}
\]
The height \(1/6\) is the exact original
section height from RJC.53; the line square
is \(-H\).
All singular line corrections have been
carried before this intersection calculation.

For the signed equation RJC.3, the actual
binding coordinate is
\[
 q_b=(t-\tau)(t/\tau)^{-\varepsilon m n_0}.
 \tag{RJC.72}
\]
It has derivative one at \(\tau\), and its
entire unit factor remains in the graph.
After retaining its constant phase, the
normal angles satisfy
\(\alpha_{\rm normal}=\phi-M\beta\), with
\(\beta=\varepsilon\gamma\).
The source-preserving boundary graph has
\(\beta=\theta+g(p)\).
K30\textquotesingle s triangular tangent shear carries
the derivatives of \(g\) and the full unit
in RJC.72; its inverse subtracts them.
It gives the physical Hopf normal sphere
with the same positive parameter orientation.
Equivalently the complete origin obeys
\[
 F_\varepsilon(v+\theta)-F_\varepsilon(v)
 =\varepsilon\theta\gamma
   -\frac{\varepsilon M}{2}\theta(\theta-1)\delta.
 \tag{RJC.73}
\]
Both the source graph and this quadratic
term are retained for each sign.

The full physical dual-line and inverse-clutch
maps of K68-\/-K69 now give, on the literal
shorter cover RJC.64-\/-68,
\[
 (k_{\rm phys},\ell_{\rm phys})
   =(3H/b,24\chi/b),\qquad
 (m_{\rm native},n_{\rm native})
   =(6\chi/b,\,3H/b-6\chi/b).
 \tag{RJC.74}
\]
These follow by the oriented degree-\(b\)
covering of integral characteristic pairs,
so no division in the finite group is used.
Substitution of the actual root Euler
numbers gives the integer table
\[
 \begin{array}{c|c|c|c}
 (a,b)&k_{\rm phys}&\ell_{\rm phys}&
 S_0=(4k_{\rm phys}-\ell_{\rm phys})/24\\ \hline
 (1,6)&kl&4kl&0\\
 (2,3)&2l(4k+1)&8l(k+1)&kl\\
 (3,2)&9k(3l+2)&12k(l+2)&k(4l+2)\\
 (6,1)&18(3l+2)(4k+1)&
 24\bigl(k(l+2)+l+1\bigr)&k(35l+22)+8l+5.
 \end{array}
 \tag{RJC.75}
\]
For each row the native pair is also exactly
\((\ell_{\rm phys}/4,
k_{\rm phys}-\ell_{\rm phys}/4)\).
The proof retains both signs of RJC.3;
the positive scalar-parameter orientation
in RJC.72-\/-73 gives these displayed pairs
for both.

\par

\medskip\noindent\textbf{RJC.11. The full fan algorithm and its exact image}

\par\nobreak

On the vertical fan \(T_j\), the primitive
return and pair are
\[
 S_j=S_0-\frac{J_0}{2}j(j-1),\quad
 k_j=k_{\rm phys}-3J_0j(j-1),\quad
 \ell_j=\ell_{\rm phys}.
 \tag{RJC.76}
\]
For every other lattice-basis fan the exact
K73 inputs, now carried by RJC.62-\/-74, are
\[
 J=J_0,\qquad \widehat H=H/b,\qquad
 \widehat\chi=\chi/b.
 \tag{RJC.77}
\]
To make the computation fully specified,
take the canonical finite word from the
actual identical-fan maps and terminating
Euclidean subtraction of K71.AF.1-\/-2.
At its vertical seed \(T_j\), start with
\[
\begin{gathered}
 V=\begin{pmatrix}
 2J-6\widehat H+6Jj(j-1)&2J(1-2j)&2J(2j-1)\\
 2J(1-2j)&2J&0\\
 2J(2j-1)&0&2J
 \end{pmatrix},\\
 D=\begin{pmatrix}
 2J+6\widehat H-24\widehat\chi-6Jj(j-1)\\
 2J(2j-1)\\2J(1-2j)
 \end{pmatrix},\\
 S=\widehat H/2-\widehat\chi-\frac J2j(j-1).
\end{gathered}
\tag{RJC.78}
\]
The matrix and vector are the actual double-cover
variation and stable-derivative invariants.
At a directed step \(P'=PQ_i\), let
\(a_*=V_{ii},b_*=D_i\), and set
\[
 \begin{split}
 B&=2J(P')^{-1}e_2e_2^tP'=(b_{uv}),\\
 L&=(b_{22}-b_{21},\,b_{11}-b_{12},\,b_{22}),\\
 (x,y)^t&=(P')^{-1}e_2,\quad
 w=(x,y,-x-y),\\
 d&=(-x-2y,-2x-y,x-y),\\
 I_r&=Jw_r+d_r(L_r-J),\quad
 W_r=Jw_r-d_r(L_r-J).
 \end{split}
 \tag{RJC.79}
\]
For the actual nonvertical step all three
target weights are nonzero. The unique
target prism solution is
\[
 \begin{split}
 C_*&=w_2^2a_*-w_2I_2,\quad
 E_*=w_2(2a_*-b_*)-W_2,\\
 V'_{rs}&=(I_r/w_r)\delta_{rs}+C_* /(w_rw_s),\\
 D'_r&=(W_r+E_*)/w_r,\qquad
 S'=S+(b_*-a_*)/4.
 \end{split}
 \tag{RJC.80}
\]
These are the full target systems of K73,
whose exterior, physical line, tangent
and two cap arrows have now been matched
to this joint source. In particular the
native resets are the maps RJC.65-\/-68,
including each full source coefficient.
Every intermediate free invariant and
increment is integral. The actual fan
return is \(S(P)\bmod28\).
The Euclidean word terminates for every
integral unimodular \(P\), and its
reciprocal chart maps give the identical
fan inverses. This proves coverage of
all basis fans, beyond the finite tests.

We compute the exact image over all fans
and all integer scalar windings for
each fixed \(k,l\).
The prism equations and their unique
solutions are linear in
\((J,\widehat H,\widehat\chi)\).
Thus scaling those three inputs scales
the full \(V,D,S\) at every step.
All four comparisons below use an actual
admissible unscaled source profile, so
its return is an integer for every fan.

For \((a,b)=(1,6)\), the inputs are
\(kl\) times the actual \(k=l=1\)
profile. Therefore every \(S(P)\)
is divisible by \(kl\).
The explicit fan
\(P=\left(\begin{smallmatrix}1&-2\\0&1\end{smallmatrix}\right)\)
has word \(Q_0Q_0\) from \(T_0\).
For the unscaled profile,
\((J,\widehat H,\widehat\chi)=(2,1/3,1/6)\);
the two selected pairs \((V_{00},D_0)\)
are \((2,2)\), then \((-4,8)\).
The seed is zero and the increments
are \(0,3\), so \(S=3kl\).
Since \(3\) is a unit modulo \(28\),
the exact image is \(kl\mathbb Z/28\).

For \((a,b)=(2,3)\), all inputs are
\(l\) times those of the actual
\((k,1)\) profile. Its admissibility
follows from
\(\gcd(3l,4k+1)=1\).
Consequently every return is divisible
by \(l\). The normalized vertical
integer is
\[
 S_j/l=k-(4k+1)j(j-1)/2.
 \tag{RJC.81}
\]
For odd \(k\), use \(j=0\) when
\(k\not\equiv0\pmod7\) and \(j=4\)
when \(k\equiv0\pmod7\).
The coefficients are \(k\) and
\(-23k-6\), both odd and nonzero
modulo \(7\) in their stated cases.
For even \(k\), use \(j=2\) when
\(k\not\equiv2\pmod7\) and \(j=3\)
when \(k\equiv2\pmod7\).
The coefficients are \(-3k-1\)
and \(-11k-3\), again odd and
nonzero modulo \(7\).
Each supplies a unit modulo \(28\);
the exact image is \(l\mathbb Z/28\).

For \((a,b)=(3,2)\), admissibility
\(\gcd(3l+2,4k)=1\) forces \(l\)
odd. All inputs are \(k\) times
those of the actual \((1,l)\)
profile; every fan return is
divisible by \(k\).
Use
\(P=\left(\begin{smallmatrix}1&-3\\0&1\end{smallmatrix}\right)\),
with word \(Q_0Q_0Q_0\) from \(T_0\).
Its full selected target-system data,
divided by \(k\), are
\[
 \begin{array}{c|c|c|c}
 \text{stage}&V_{00}/k&D_0/k&(D_0-V_{00})/(4k)\\ \hline
 T_0&-42l-28&54l+20&24l+12\\
 Q_0&-204l-136&312l+176&129l+78\\
 Q_0Q_0&-474l-316&1110l+692&396l+252.
 \end{array}
 \tag{RJC.82}
\]
Adding these increments to the seed
\(4l+2\) gives \(553l+344\).
This is odd for odd \(l\), and
\(553l+344=7(79l+49)+1\).
It is a unit modulo \(28\).
The exact image is \(k\mathbb Z/28\).

For \((a,b)=(6,1)\), put
\(A=k(35l+22)+8l+5\) and
\(J=(3l+2)(4k+1)\).
The four explicit vertical integers
needed are
\[
 \begin{split}
 S_0&=A,\\
 S_4&=-k(37l+26)-10l-7,\\
 S_2&=k(23l+14)+5l+3,\\
 S_3&=-k(l+2)-l-1.
 \end{split}
 \tag{RJC.83}
\]
When \(kl\) is even, \(S_0,S_4\)
are both odd. Their residues modulo
\(7\) are \(k+l+5\) and
\(-k(2l+5)-3l\).
Common zero would give \(k=2-l\)
and \(2(l-4)^2=0\pmod7\), hence
\((l,k)=(4,5)\pmod7\).
When \(kl\) is odd, \(S_2,S_3\)
are both odd. Their common zero
has \(l\ne5\pmod7\), since
\(S_3\) is then nonzero.
Eliminating \(k\) gives
\(3l^2+11l+6=3(l-4)^2=0\pmod7\),
and again \((l,k)=(4,5)\pmod7\).
But that residue makes both
\(3l+2\) and \(4k+1\) divisible
by \(7\), contradicting the
actual coprimality in RJC.3.
Thus one of the displayed two
odd integers in the relevant
parity case is a unit modulo
\(28\). The exact image is all
of \(\mathbb Z/28\).

The result is the full table
\[
 \begin{array}{c|c|c}
 (a,b)&\text{holomorphic winding group}&
 \text{all-fan, all-loop return image}\\ \hline
 (1,6)&2kl\,\mathbb Z&kl\,\mathbb Z/28\\
 (2,3)&l(4k+1)\,\mathbb Z&l\,\mathbb Z/28\\
 (3,2)&2k(3l+2)\,\mathbb Z&k\,\mathbb Z/28\\
 (6,1)&(3l+2)(4k+1)\,\mathbb Z&\mathbb Z/28 .
 \end{array}
 \tag{RJC.84}
\]
The displayed nonvertical witnesses use
the entire prism solves, not a polynomial
extrapolation in the fan entries.

\par

\medskip\noindent\textbf{RJC.12. Every holomorphic scalar path and both mapping-torus orientations}

\par\nobreak

Let a holomorphic scalar path on its
parameter annulus have a holomorphic
endpoint identification on these fixed
marked objects. RJC.5-\/-8 forces equal
scalar values at its ends and winding
\(J=sJ_0\), \(s\in\mathbb Z\).
On the universal logarithmic strip
choose its full logarithm \(L\), with
\(L(\theta+1)-L(\theta)=2\pi iJ\).
The explicit deformation
\[
 L_u(\theta)=(1-u)L(\theta)
       +u\bigl(L(0)+2\pi iJ\theta\bigr),
 \qquad 0\le u\le1
 \tag{RJC.85}
\]
is holomorphic in the annular parameter
and preserves the integral winding.
The contraction domain is the original
punctured sufficiently small scalar
disc; convexity of its logarithmic
real bound keeps this deformation
inside it. On the compact deformation
interval the original source and line
norm bounds give common strictly
contracting cap tubes.

Use \(L_u\) in the full equations
RJC.33-\/-41, keeping every integer
period, finite character and reference
value. Use RJC.64-\/-68 on every selected
chart. Their inverse arrows have the
same continuous dependence. These
maps form the actual smooth proper
family over the deformation square.
Local lifts of its parameter vector
field, patched by a partition of
unity and chosen to match the given
end collars, integrate for the
whole compact interval. This gives
the relative endpoint isotopy and
preserves the specified source,
line and cap comparisons.

The ratio of its closing to the
chronological \(s\)-th power of
RJC.15 is a physical scalar by
RJC.56-\/-58. Choosing any logarithm
of that scalar, the actual path
of multiplying by \(\exp(u\log c)\)
extends on every point line,
cap and fan and gives its isotopy
to identity. Thus all holomorphic
closings at that winding have
the same smooth return.
Every integer \(s\) is attained
on an annulus about \(w_*\ne0\)
by the explicit scalar parameter
map
\[
 w\longmapsto
 w_*\left(\frac{w}{w_*}\right)^s.
 \tag{RJC.86}
\]
Pulling back the actual shorter
family gives its closing. Negative
\(s\) uses the actual inverse
maps; \(s=0\) is the constant
family. At a fixed fan \(P\) the
complete set of returns is
\[
 [\Phi]=s\,S(P)\pmod{28}.
 \tag{RJC.87}
\]
Combining this with the unit
witnesses proves the images
in RJC.84 for all paths and
all holomorphic closings.

Finally use the complete seam
and disk-transfer maps of the
proved mapping-torus construction.
For the positive seam
\((x,1)\sim(\Phi(x),0)\),
\[
 T_+(\Phi)\cong
 (S^1\times S^6)\mathbin{\#}\Sigma_{sS(P)}.
 \tag{RJC.88}
\]
The inverse seam has
\(\Sigma_{-sS(P)}\).
The physical inverse-clutch
identification is
\(\lambda_{\rm physical}(F)
=-\Sigma_{F^{-1}}=\Sigma_F\);
it is the one transported in
RJC.72-\/-74.
For each of the two spin
structures the corresponding
Eells-\/-Kuiper values are
\[
 \mu(T_+)=\frac{sS(P)}{28},
 \qquad
 \mu(T_-)=-\frac{sS(P)}{28}
 \quad\text{in }\mathbb Q/\mathbb Z.
 \tag{RJC.89}
\]
Every nonzero residue obstructs
an orientation-preserving product
identification in that orientation;
reversing the seam changes its
sign by the explicit inverse.

\par

\medskip\noindent\textbf{RJC.13. Exact checks and remaining scope}

\par\nobreak

The primary executable reconstructs
the original odd-line Miller
pullbacks, inverse-source
coefficients, root characters,
signed norm equations, both
root exponents and the full
cyclic scalar identity.
It calculates the four page
pairs and performs fresh
nine-unknown target-prism solves,
then checks the actual unit-fan
witnesses and the unbounded
modulo-seven identities.
The geometric statements above
use the complete source, line,
root, reference and parameter
maps before their numerical
characteristic evaluations.
The bounded tests support those
proofs; the terminating fan
construction and the exact
residue arguments give their
unbounded conclusions.

This proof has the two actual
additive caps of RJC.2 and
signed equation RJC.3. The
additional ordinary norm is
zero throughout. The prescribed
ordinary references, all finite
characters and every displayed
root order are included.
Further ordinary clutch slots
with nonzero total norm require
their actual redistributed
cochain comparison before they
enter this statement. Other
marked source-moduli loops are
also outside this scalar-path
classification. Part3.4 remains
incomplete. Independent full
geometric audit and export are
still pending for this
continuation.

% End expanded verified cp36 input


\par\bigskip\noindent\textbf{Complete supporting proof 2: JFP.}\par


% Expanded verified cp36 input: notes/part3_e6_iii_joint_endpoints_jfp_proof.tex
\par

\medskip\noindent\textbf{Independent full comparison for the original E6 source with both additive caps}

\par\nobreak

This proof starts from the actual F46 source and Cartier construction,
the K67 shorter parameter family, the K73 native products, and the
completed K114--115 endpoint and scalar-family maps. The root exponents,
cyclic endpoint, characteristic inputs and fresh nine-variable prism
systems were derived and executed before the primary RJC manuscript
was read. The present proof also records the subsequent complete
RJC.1--13 geometric review. Every point below uses the original source
and its specified references. The ordinary extra norm is zero.

\par

\medskip\noindent\textbf{JFP.1. Source, both root maps, and all point divisors}

\par\nobreak

Keep
\[
 y^2=x^3+tx+t^2,\quad G=(0,t),\quad q=t+4/27,
 \quad \mathcal M_a=\mathcal O([aG]-[O]),\quad ab=6.
\tag{JFP.1}
\]
The four domains, with their original positive orders, are
\[
\begin{array}{c|c|c|c}
(a,b)&m=m_\infty&h=m_0&\text{root fibres at infinity and zero}\\\hline
(1,6)&3l, l\ge1&4k, k\ge1&\text{smooth, smooth}\\
(2,3)&3l, l\ge1&4k+1, k\ge0&\text{smooth, III}\\
(3,2)&3l+2, l\ge0&4k, k\ge1&\text{IV, smooth}\\
(6,1)&3l+2, l\ge0&4k+1, k\ge0&\text{IV, III}.
\end{array}
\tag{JFP.2}
\]
The prescribed signed norms obey
\[
 h n_\infty+m n_0=\varepsilon,\qquad \varepsilon=\pm1.
\tag{JFP.3}
\]
In particular \(\gcd(m,h)=1\) and both local norms are primitive.
This is F46\textquotesingle s integral sphere equation. Its actual source/cap
chain maps have degree-two map an integral isomorphism and
degree-one and degree-three determinants \(-\varepsilon\).
Their fundamental-group presentation is cyclic of order one.
The common blowup and its two ruling projections carry these maps
through each prescribed lattice-basis fan, preserving the primitive
normal circle and the physical line. Thus the F46 sphere and its
vanishing \(H^2,H^4\) apply on every fan used here.

At infinity the full substitutions are
\[
\begin{array}{c|c|c}
m&(t,x,y)&\text{equation}\\\hline
3l&(z^{-3l},z^{-2l}X,z^{-3l}Y)&Y^2=X^3+z^lX+1\\
3l+2&(z^{-m},z^{-2(l+1)}X,z^{-3(l+1)}Y)&
Y^2=X^3+z^{l+2}X+z^2.
\end{array}
\tag{JFP.4}
\]
At zero they are
\[
\begin{array}{c|c|c}
h&(t,x,y)&\text{equation}\\\hline
4k&(z^{4k},z^{2k}X,z^{3k}Y)&Y^2=X^3+X+z^{2k}\\
4k+1&(z^h,z^{2k}X,z^{3k}Y)&Y^2=X^3+zX+z^{2k+2}.
\end{array}
\tag{JFP.5}
\]
The deck multiplies \(z\) by its specified root of unity and
\(X,Y\) by the inverse displayed weights. Each original coordinate
is therefore invariant. The reciprocal powers give the inverse
on the original annulus. The reduced equations retain their A1 or
A2 resolution charts.

The full good-cubic point-line graph uses the quotient weights
\((2,1)\). On its common regular refinement the rays
\((1,0),(2,1),(1,1),(1,2),(0,1)\) give the discrepancy
\(\min(v_\eta,2v_z)\). After the actual pullback \(w=z^l\),
the correspondence is
\[
 \beta^*\mathcal O(aG'-O)=\pi^*\mathcal M_a+D_{aG'}-D_O,
 \qquad D_P(v)=\min(v_\eta,2l v_z).
\tag{JFP.6}
\]
Each cone map is its two original monomials in \(\eta,z\), with
the integral inverse on the common torus. The strict point-section
factor is retained in that generator. This carries both the original
point module and all its exceptional factors.

On the good quartic root, the actual action
\((w,X,Y)\mapsto(iw,-X,iY)\) fixes \(O,T=(0,0)\) and exchanges
\((i,0),(-i,0)\). Blow up those four points and take the quotient.
The exceptional images have squares \(-4,-4,-2\), the central
curve has square \(-1\), and the fibre is
\(4C+D_O+D_T+2D_2\). Contracting \(C\), then the resulting
\((-1)\)-curve \(D_2\), gives the original two \((-2)\)-curves
with intersection scheme of length two. The invariant maps (JFP.5)
give their actual tangency. The marked sections avoid these two
contractions. The line correspondence is
\[
 \beta^*\mathcal O(aG'-O)
 =\pi^*\alpha^*\mathcal M_a+E_{aG'}-E_O.
\tag{JFP.7}
\]
For \(w=z^k\), the section graphs resolve the ideals
\((\eta,z^k)\). The complete successive cone maps are
\[
 \eta=u^jv^{j+1},\quad z=uv\quad(0\le j<k),
 \qquad \eta=uv^k,\quad z=v
\tag{JFP.8}
\]
on the last cone. Their determinants are \(-1\). They pull the
ideal to the indicated generator with valuation
\(\min(v_\eta,k v_z)\); the corresponding reciprocal module maps
are the inverse. For odd \(a\), translation exchanges the two
marked exceptional contributions. Its actual Hom coefficient
carries that exchange, as computed in JFP.2 below.

For the reduced III root the original formal parameter is
\(\tau=-x/y\) and \(\tau'=z^k\tau\). For \(a=2,6\),
\(\tau'_{aG'}=z^kA_a(z^h)\), where \(A_a(0)=a\).
The two actual point/O ratios are
\[
 \frac{p-A_a(z^h)}p,\quad 1-vA_a(z^h),
 \qquad \tau'=z^kp,\quad z^k=\tau'v.
\tag{JFP.9}
\]
The intermediate charts are (JFP.8). The A1 chart is
\(X=z\alpha,Y=z\beta\),
\(\beta^2=z\alpha^3+\alpha+z^{2k}\). The slopes
\((Y\mp z^{k+1})/X\) give the two inverse source translations
on their resolved graphs. For the reduced IV root, use the full
A2 charts, beginning with \(X=zX_1,Y=zY_1\), whose two outer
component labels are \(Y_1=\pm1\). Its point graph is
\(\tau'/(\tau'-z^{2l+1}L_a(z))\), with the actual functions
\[
 L_3(z)=-\frac{64+8z^m}{512+96z^m+3z^{2m}},
 \quad L_3(0)=-\frac18,
 \qquad L_6(0)=-\frac14.
\tag{JFP.10}
\]
Here \(L_a=-z^{-2l-1}X_{aG'}/Y_{aG'}\), so its complete
rational function, and its reciprocal module map, stay in the
correspondence. The distinct constants are obtained by the original
\(3G\) and doubled \(6G\) formulae. This was the one source-graph
correction found in the full primary read; the preceding primary
and independent derivation are preserved in their attempt records.

At either reduced root the point \(aG'\) lies in the identity
component. The original pulled-back line and the displayed root
line have zero degrees on all components. The blowup Picard
decomposition descends the original line through every contracted
curve. Their remaining vertical divisor is an integer whole fibre,
because the resolved Kodaira intersection matrix has primitive
kernel its fibre vector. Division by the actual base equation to
that integer power gives the line map; multiplication gives its
inverse. The power and its deck character remain specified.

\par

\medskip\noindent\textbf{JFP.2. Original Miller maps and the odd-line coefficients}

\par\nobreak

Let \(T(p)=p-G\). On the original source retain
\[
\begin{split}
 f_1&=1,\qquad f_2=\frac{y-t-x/2}{x-1/4},\\
 f_3&=\frac{(x/4)(y-t-x/2)}{(y+t+x/2)/4-(1/4+2t)x},\\
 (x_3,y_3)&=(8t+64t^2,3t+96t^2+512t^3),\\
 L_3^{\rm tan}&=(3x_3^2+t)/(2y_3),\quad
 x_6=(L_3^{\rm tan})^2-2x_3,\\
 y_6&=L_3^{\rm tan}(x_3-x_6)-y_3,\quad
 f_6=f_3^2\frac{y-y_3-L_3^{\rm tan}(x-x_3)}{x-x_6}.
\end{split}
\tag{JFP.11}
\]
Their divisors are \(a(G)-(aG)-(a-1)(O)\). The full one-step
negative lift and its positive inverse therefore are
\[
 B_{a,1}(p)=f_2(p)^a\frac{f_a(Tp)}{f_a(p)},\qquad
 U(p,v)=\bigl(p+G,B_{a,1}(p+G)^{-1}v\bigr).
\tag{JFP.12}
\]
The inverse coefficient is evaluated at the inverse source point.
These are Hom-module identities even where a displayed rational
coefficient has a zero or pole; the point/O Cartier generators
carry those factors.

On the two original III components use
\(x=p^{-2},y=-p^{-3}\) and
\(x=t\alpha,y=t\beta\),
\(\alpha+t\alpha^3=\beta^2-1\). The directed source maps are
\(\beta'=p-1,p'=\beta-1\). Substitution into (JFP.11) gives
\(f_2,f_3,f_6\) component orders \((0,1),(0,1),(0,3)\).
Consequently the one-step Hom orders for \(a=1,2,3,6\) are
\((0,1),(1,1),(1,2),(3,3)\). The chronological \(b\) factors
have order three on both components. At every other finite fibre
the two translated point factors have the same generic order;
\(f_2\) has no whole-fibre zero. The actual original translation
Hom has class \(F\), so its infinity coefficient is \(-2\).
The specified suspension section is exactly
\[
 B_0=C_a\frac q{t^3}U^{-b},\qquad
 \rho=\frac{B_0^{-1}}{U^b}=\frac{t^3}{C_aq}.
\tag{JFP.13}
\]
The ratio of the two sections has zero divisor on the compact
regular source and is an invertible constant. Its evaluation at
the prescribed reference fixes the actual \(C_a\in\mathbb C^*\).

The new odd-line good-root coefficients can also be computed
without relying on a divisor alone. At the finite good root,
\(w=z^k\), direct substitution gives
\[
 f_{2,\rm old}=w^2\frac{wY-w^2-X/2}{w^2X-1/4},\quad
 f_{3,\rm old}=w^2\frac{X(2wY-2w^2-X)}
 {2wY+2w^2-X-16w^4X}.
\tag{JFP.14}
\]
The central source map is \((X,Y)\mapsto(1/X,-Y/X^2)\).
The full root functions specialize to \(f_1=1,f_2=2X,
f_3=X,f_6=6X^3\). Hence the one-step negative coefficients
are \(2^aX\) for odd \(a\), and \(2^a\) for even \(a\).
The positive inverse coefficients are \(X/2^a\) and \(1/2^a\),
respectively. For the two odd rows their chronological \(b\)-fold
positive norm is \(1/64\). The exchanged \(E_{aG'}-E_O\)
generators in (JFP.7) are carried by these exact coefficients.

At the good infinity root, \(w=z^l\),
\[
 f_{2,\rm old}=w^{-1}\frac{Y-1-wX/2}{X-w^2/4}.
\tag{JFP.15}
\]
On \(Y^2=X^3+1\), the mutually inverse source maps are
\[
 (X,Y)\mapsto\left(-\frac{2X}{Y+1},\frac{Y-3}{Y+1}\right),
 \qquad
 (X,Y)\mapsto\left(\frac{2X}{Y-1},\frac{-Y-3}{Y-1}\right).
\tag{JFP.16}
\]
The positive \(M_1\) coefficient is \(X/2\), whose three
chronological factors give \(-1/2\) and whose six give \(1/4\).
The positive \(M_2\) coefficient is \((Y+1)/(2X)\), whose
three factors give \(1/4\). The source third iterate is the
identity. These full source and line calculations carry the
good-root point corrections (JFP.6); their inverse maps retain
the full off-central functions (JFP.14--15).

\par

\medskip\noindent\textbf{JFP.3. Both root lifts, division lifts, and finite characters}

\par\nobreak

The extending positive root lift \(L_i\) is chosen with its
specified rigidification and \(L_i^b=B_0^{-1}\). The entire
ordinary-to-root comparison is
\[
 U=z^{e_i}u_i L_i,\qquad
 \begin{array}{c|c|c}
 i&e_i&u_i^b\\\hline
 0&-3h/b&C_aq\\
 \infty&2m/b&C_a(1+4z^m/27).
 \end{array}
\tag{JFP.17}
\]
To construct it, apply the exact point modules of JFP.1 to the
root chord lift. Its horizontal divisor is the prescribed Hom
divisor. Its remaining vertical orders have zero component degrees;
the primitive fibre kernel makes them one integer base power.
Remove that displayed power. The section and its inverse then
have zero orders at every codimension-one point. Normal local
rings are intersections of their height-one valuation rings, so
both extend through component and section intersections. The
nonconstant logarithmic deck commutator is removed by finite
averaging. For a norm-zero logarithm \(g\), the explicit average
\(-m_i^{-1}\sum_{j=0}^{m_i-1}jD_i^jg\) gives the required
coefficient correction. The invariant analytic \(b\)-th root
of the norm ratio, with the prescribed reference value, fixes
\(L_i^b=B_0^{-1}\). Its reciprocal gives the inverse correction.
The retained constant character is obtained from (JFP.17).

The coefficient convention matters. Conjugating multiplication
by \(z^e\) by \(D:z\mapsto\zeta z\) evaluates that function
at \(z/\zeta\), giving \(\zeta^{-e}z^e\). Original \(U\)
is invariant, so
\[
 D_iL_iD_i^{-1}=\zeta_{m_i}^{e_i}L_i,\qquad
 \chi_\infty=\zeta_b^2,\quad\chi_0=\zeta_b^{-3}.
\tag{JFP.18}
\]
Thus the four character pairs are
\((\zeta_3,-1),(\zeta_3^2,1),(1,-1),(1,1)\).
The norm equation (JFP.3) gives
\(\chi_i^{-\varepsilon}=\zeta_b^{Mn_i/m_i}\),
with \(M=mh\). Every prescribed common finite line character
cancels in this conjugation, while it remains in the cap operator.

For completeness the cap division lift has character one.
The marked division section \(R_i\) satisfies \(m_iR_i=bG'\).
At a good root it is the unique analytic inverse branch with
\(R_i(0)=O\). At a reduced root multiplication by \(m_i\)
is invertible on the additive identity group and on its finite
component group, so it gives the unique marked analytic section.
The deck preserves it. For \(P_i=aG'\), the full chord ratio
\[
 g_{P_i,R_i}(p)=
 \frac{\ell_{P_i,R_i}(p)}{x(p)-x(P_i+R_i)}
\tag{JFP.19}
\]
has divisor \((P_i)+(R_i)-(P_i+R_i)-(O)\) and is the Hom
coefficient for \(p\mapsto p-R_i\). Its whole-fibre power is
removed through the actual generators (JFP.6--10). At a good root
the central source translation is the identity and the central
line map is a constant. At a reduced root of order greater than one, the actual \(aG'\)
section limits to \(O\), so its point/O line is trivial with
the displayed identification; the central translation lift is
again a constant in that identification. At the finite order-one
root, the full graph is \(\tau'=A_a(t)\) with
\(A_a(0)=a\), as in (JFP.9), and the deck is the identity.
Its commutator is therefore one directly on that actual point line.
The commutator\textquotesingle s
constant character is therefore one. Finite averaging and the
invariant norm root give the complete commuting operator
\(\widetilde F_i^{m_i}=B_0\).

With its retained character \(d_i\), the actual cap generator is
\[
 A_{i,\lambda}=R_{\lambda,i}^{n_i}\zeta_{m_i}^{d_i}
 \widetilde F_i^{n_i}D_i,\qquad
 R_{\lambda,i}^{m_i}=\lambda,\qquad
 A_{i,\lambda}^{m_i}=B_\lambda^{n_i}.
\tag{JFP.20}
\]
All factors commute with the suspension. For integers \(u_i,v_i\)
with \(n_iu_i+m_iv_i=1\), the primitive generator
\(A_i^{u_i}B^{v_i}\) is the strict contraction with its actual
scalar root. This proves the cap shell and its physical-circle
freeness on the stated domain for either sign, with the same
Cartier and finite-character factors.

\par

\medskip\noindent\textbf{JFP.4. Complete annular cochains and their integer restoration}

\par\nobreak

Set \(q_0=t,q_\infty=1/t\), \(m_0=h,m_\infty=m\), and retain
\[
 s_{i\alpha}=\log_\alpha q_i/(2\pi i),\quad
 z_{i\alpha}=e^{2\pi i s_{i\alpha}/m_i},\quad
 c_{\lambda,i}=n_i\log R_{\lambda,i}/(2\pi i)+d_i/m_i,
 \quad f_{i\alpha}=\operatorname{Log}_\alpha\widetilde F_i.
\tag{JFP.21}
\]
The original-coordinate and point-module root map is \(I_z\).
The complete gluing is
\[
 \Phi_{i\alpha}=\operatorname{Exp}
 (n_i s_{i\alpha}f_{i\alpha}+s_{i\alpha}c_{\lambda,i}\delta
       +h_{i\alpha})I_{z_{i\alpha}}.
\tag{JFP.22}
\]
Its source translation is \(-n_i s_{i\alpha}R_i\). On an
overlap labelled \(k\),
\(f_{i\beta}=D_{i*}^kf_{i\alpha}+\Pi_{i\alpha\beta}\), and
the retained reference cochain satisfies
\[
 h_{i\beta}-D_{i*}^kh_{i\alpha}
 =-n_i s_{i\alpha}\Pi_{i\alpha\beta}
   +j_{i\alpha\beta}\operatorname{Log}_\beta B_\lambda
   +\Lambda_{i\alpha\beta}.
\tag{JFP.23}
\]
The periods \(\Pi,\Lambda\), suspension integers \(j\), and
reference value are unchanged. Starting from its prescribed smooth
extension \(h^0\), subtract the residual with
\(k_\alpha=-\sum_\beta\chi_\beta E_{\alpha\beta}\), then
subtract the common Cauchy integral
\(\pi^{-1}\int\bar\partial k(w)/(q-w)\,dA(w)\).
The residual has compact subannular support. Fixing the common
holomorphic constant at the original reference gives exactly that
reference value. Subtracting the target-sheet logarithm proves
\[
 \Phi_{i\beta}=A_{i,\lambda}^{k}B_\lambda^{j_{i\alpha\beta}}
 \Phi_{i\alpha},\qquad \Phi_{i\alpha}B_\lambda=B_\lambda\Phi_{i\alpha}.
\tag{JFP.24}
\]
The inverse uses the negative Lie translation at the inverse
source point, followed by the inverse root map and reciprocal
Cartier factors.

For \(\lambda_\theta=e^{2\pi i\theta}\lambda_*\), retain the
actual integers \(J_{i\beta}-J_{i\alpha}=j_{i\alpha\beta}\).
The full restoration formula is
\[
 \Phi_{\theta,i\alpha}
 =e^{2\pi i n_i\theta s_{i\alpha}/m_i}
 B_\theta^{J_{i\alpha}}
 (B_*^{-J_{i\alpha}}\Phi_{*,i\alpha}).
\tag{JFP.25}
\]
At integer winding \(J\), the last two factors have scalar ratio
\(e^{2\pi iJ J_{i\alpha}}=1\), and the first gives
\(\Phi_{\theta+J,i\alpha}=z_{i\alpha}^{n_iJ}\Phi_{\theta,i\alpha}\).
This establishes both root exponents used below through the full
reference, with every integral period still present.

\par

\medskip\noindent\textbf{JFP.5. Primitive endpoint and its exact cyclic phase}

\par\nobreak

Put \(J_0=M/b\) and retain the actual one-gamma forward closing
\(h_0=c_0q^\varepsilon t^{-mn_0}\). The single-valued integer
\[
 d_0=(3\varepsilon-mn_0)/b=(2\varepsilon+hn_\infty)/b
\tag{JFP.26}
\]
has respectively the integral expressions
\(2kn_\infty+ln_0\),
\((4k+1)n_\infty+2ln_0\),
\(6kn_\infty+(3l+2)n_0\), and
\(3(4k+1)n_\infty+2(3l+2)n_0\).
Choose the specified root \(\gamma_0^b=c_0C_a^{-\varepsilon}\).
The full positive endpoint over \(\lambda=r^M\) is
\[
 E_r=\gamma_0r^{-\varepsilon J_0}t^{d_0}U^\varepsilon,
 \qquad r\longmapsto\zeta_b r.
\tag{JFP.27}
\]
At zero the full base exponent is
\(hd_0-3h\varepsilon/b=-J_0n_0\); at infinity it is
\(-md_0+2m\varepsilon/b=-J_0n_\infty\).
The reference factors of (JFP.25) cancel exactly these powers.
The remaining cap coefficient is
\(\gamma_0r^{-\varepsilon J_0}u_i^\varepsilon L_i^\varepsilon\),
and (JFP.18) gives its exact finite-generator square. The inverse
is the reciprocal coefficient at the inverse source point.

The chronological product keeps the scalar phase
\(\exp[-\pi i\varepsilon J_0(b-1)]\). In all four rows
\(J_0(b-1)\) is even. With the full constants in (JFP.13),
\[
 E_{\zeta_b^{b-1}r}\cdots E_{\zeta_b r}E_r
 =h_0 B_\lambda^{-\varepsilon}.
\tag{JFP.28}
\]
The K67 bounded origin uses
\(c=\exp[\pi i\varepsilon J_0(b-1)/b]\). The exact family
intertwiner is physical multiplication by
\[
 F_r=r^N,\qquad N=\varepsilon J_0(b-1)/2\in\mathbb Z,
 \qquad F_{\zeta_b r}E_r=(cE_r)F_r.
\tag{JFP.29}
\]
The four exponents are \(5\varepsilon kl\),
\(\varepsilon l(4k+1)\), \(\varepsilon k(3l+2)\), and zero.
The phase has this explicit map and reciprocal inverse; it remains
in the complete parameter family.

\par

\medskip\noindent\textbf{JFP.6. Intrinsic elliptic quotient and base rigidity}

\par\nobreak

The F46 topology in JFP.1 gives \(H^2=H^4=0\). A horizontal
polar divisor of a meromorphic function restricts to a nonempty
effective divisor on a compact regular two-torus. Pairing with a
flat Kähler form is positive, whereas its restricted cohomology
class is zero. Thus all poles are vertical. The function is
constant on each regular compact torus. At a special component,
the full transverse map \(u\mapsto t_0+u^e\) and invariance
force Laurent powers divisible by \(e\). Consequently pullback
gives the precise field isomorphism \(\mathbb C(t)\cong\mathbb C(X)\),
whose inverse is fibre restriction followed by this Laurent descent.

The full generic period module is
\[
\begin{gathered}
 \gamma=(-bg,\ell),\quad e_1=(1,0),\quad e_2=(\tau,ag),
 \quad\delta=(0,1),\\
 M_\alpha=
 \begin{pmatrix}1&0&0\\bv_\alpha&A_\alpha&0\\
 c_\alpha&a v_\alpha^tJ A_\alpha&1\end{pmatrix},
 \quad J=\begin{pmatrix}0&-1\\1&0\end{pmatrix}.
\end{gathered}
\tag{JFP.30}
\]
The full theta and Miller constant is in \(\ell\). Multiplication
gives the retained central cochain
\(c_{\alpha\beta}=c_\alpha+c_\beta+abv_\alpha^tJ A_\alpha v_\beta\).
The Abel cocycle is nonzero over \(\mathbb Q\). A rational
coboundary, after clearing denominators and component orders,
would supply a global Lie logarithm of a nonzero multiple of
\(G\). At the minimal quartic root divide the logarithm of
\(2G'\) by its actual factor \(w\); at the cubic root divide
the logarithm of \(3G'\) by \(w^2\). The latter has leading
coefficient \(-1/8\). Both finite elliptic actions have no
invariant period; at the I1 the invariant annular period extends
through \(xy=q\). The original Lie transition is
\(u_\infty(q_\infty)=q_\infty^{-1}u_0(q_\infty^{-1})\).
Its entire global coefficient is zero. But original \(3G\)
has infinity value \((64,512)\), with nonzero additive parameter
\(-1/8\), and every nonzero multiple remains nonzero. The
contradiction proves the Abel assertion. The higher maps JFP.1
are the actual pullbacks of these minimal diagrams.

The vector and dual equations of (JFP.30) now give exactly
\(\Lambda^{\pi_1}=\mathbb Z\delta\) and
\((\Lambda^*)^{\pi_1}=\mathbb Z\psi\), where \(\psi(\gamma)=1\)
and it vanishes on the other three generators. Both are primitive.
The kernel cover is the actual nonzero point line with its Cartier
modules. The derivative of a lifted endpoint is periodic for the
compact torus lattice and hence constant. The full affine map and
inverse are
\[
\begin{gathered}
 (z,w)\mapsto(Az+z_0,Cz+e w+w_0),\\
 (z',w')\mapsto((z'-z_0)/A,\,
 e[w'-w_0-C(z'-z_0)/A]),\\
 A\ne0,\qquad e=\pm1.
\end{gathered}
\tag{JFP.31}
\]
On the kernel lattice its matrix is
\(K=\left(\begin{smallmatrix}p&q_1&0\\r&s&0\\u&v&e\end{smallmatrix}\right)\),
with \(ps-q_1r=1\). Its inverse has last row
\((e(vr-us),e(uq_1-vp),e)\) and upper block
\(\left(\begin{smallmatrix}s&-q_1\\-r&p\end{smallmatrix}\right)\).
The positive determinant is the complex area identity. The first
coordinates therefore descend to inverse elliptic source maps.
The original cubic gives \(j_0(t)=6912/(27t+4)\). Equality under
that actual elliptic isomorphism forces the base Möbius map to be
the identity.

The order-three infinity block and the selected nonzero unipotent
have simultaneous rational commutant \(\mathbb QI\): a nonscalar
element of \(\mathbb Q(\zeta_3)\) generates a field, which has no
nonzero nilpotent. The off-diagonal equations and nonzero Abel
class leave the full lattice commutant \(eI+c\delta\psi\).
Complex linearity on the independent vectors \(e_1,\delta\)
sets the derivative to \(eI\). The remaining suspension equation
makes the scalar logarithm difference an integer. The two ends
therefore have equal original scalar values.

\par

\medskip\noindent\textbf{JFP.7. Global section extension and all holomorphic endpoints}

\par\nobreak

Pull an endpoint and its inverse to the normalized root fibre
product \(z_i^{m_i}=q_i\). A nonzero line vector over the
O-section kills the root meridian; the remaining vanishing cycles
lie in \(\ker\psi\). The actual graphs in JFP.1 supply this
lift on the whole root source. Its commutator with the cap
generator is \(B^s\). Iterating the relation \(m_i\) times gives
\(B^{m_is}=1\); strict contraction implies \(s=0\).
The source displacement is an equivariant holomorphic root
section. Descend by (JFP.4--5) and resolve its graph in the proper
minimal original source.

At the selected source a coefficient pole of order \(-K\)
is cancelled by the actual \(B^K\) reindex. It shifts the source
displacement by \(-bKG\). The remaining two character factors
are holomorphic units; their complete vertex arrows are written
in JFP.8. Their inverses extend over both vertices and the open
axes. After this selected reindex the joining torsor has
presentation
\[
 \langle x,y\mid x^m=y^h=xy=1\rangle.
\tag{JFP.32}
\]
It is cyclic of order \(\gcd(m,h)=1\). Thus the integer joins
are globally trivial, with the one constant join fixed at the
prescribed reference. This proves extension to a global source
section.

The source is rational by the inverse birational chart
\[
 u=(y-t-x/2)/x,\quad v=(y+t+x/2)/x,\quad
 x=uv+1/4,\quad t=x(v-u-1)/2,\quad y=x(u+v)/2.
\tag{JFP.33}
\]
Its ten classes \((O,F,C_0,\ldots,C_5,A_1,G)\) have the
negative E6 block with edges \((0,1),\allowbreak(1,2),\allowbreak(2,3),\allowbreak(3,4),\allowbreak(2,5)\),
\(A_1^2=-2\), \(O^2=G^2=-1\), \(OF=GF=1\), and
\(GC_0=GA_1=1\). All other unspecified intersections are zero.
The complete determinant is \(-1\). These classes span the
rank-ten unimodular lattice; quotienting the first nine is free
cyclic generated by \(G\). Hence the section group is exactly
\(\mathbb ZG\), with height \(2-4/3-1/2=1/6\).

A positive endpoint is therefore the full map \(H(t)U^j\).
For winding \(J\), the two root comparisons and selected reindex
give the actual divisor orders
\[
 \operatorname{ord}_0H=3j/b-n_0J/h,\quad
 \operatorname{ord}_\infty H=-2j/b-n_\infty J/m,\quad
 \operatorname{ord}_{t_c}H=-K.
\tag{JFP.34}
\]
Elsewhere the Hom and its inverse have no zero. Summing the
principal divisor gives
\(j=bK+b\varepsilon J/M\), so \(J\in J_0\mathbb Z\).
The endpoint (JFP.27), its chronological powers and its inverse
attain every such integer.

For one object \(J=0\), so \(j=bK\),
\(H=c(t^3/q)^K\), and the whole map is
\(cC_a^K\lambda^KB_\lambda^{-K}\). Its quotient is the
displayed physical scalar. Faithfulness follows because a
nonzero suspension deck translates by a nonzero multiple of
the non-torsion \(bG\). A negative automorphism has sixteen
transverse fixed points on each regular two-torus. Its compact
analytic equalizer has a horizontal curve, whose positive base
degree contradicts pairing with \(\pi^*H^2(\mathbb P^1)=0\).
Thus \(\operatorname{Aut}_{\rm hol}(X)=\mathbb C^*_{\rm physical}\),
with the inverse identification just given. Two positive
endpoints at the same winding differ by precisely that scalar.

The full selected negative arrow is
\((Z,V,q)\mapsto(Z^{-1},V^{-1},q)\). Opposite triangles are
translated by \(-p-u-v\); the vertex exchange retains all
\(\kappa g,\alpha\eta\) factors and their reciprocals, and
extends on every fan. The source/root lifts of a negative
endpoint exist by the same point-vector and inverse-map
argument. Compose at the selected charts with this inversion
to remove its selected integer join. Both lifts satisfy
\(\widetilde F B=B^{-1}\widetilde F\). On the infinity annulus
the full reference word is \(A_\infty B^\nu\). One root circle
has \(B^N\), \(N=n_\infty+m\nu\). After fixing one constant
deck on a cut annulus, continuation yields
\[
 B^N\widetilde F B^{-N}=B^{2N}\widetilde F.
\tag{JFP.35}
\]
Single-valuedness and contraction force \(N=0\). But
\(\gcd(m,n_\infty)=1\) and every allowed \(m\ge2\), so
\(n_\infty+m\nu\ne0\). This proves the negative exclusion
through the actual root word, including the quadratic reduced
case, and completes the exact winding group classification.

\par

\medskip\noindent\textbf{JFP.8. Complete parameter comparison and the finite cover}

\par\nobreak

Let \(v=\log r/(2\pi i)\), and keep the bounded source period
\(\Gamma\), ordinary logarithm \(H_0=\varepsilon s+U_0(q)\),
and outward suspension \(\gamma=(\Gamma,C-s-Mv)\).
The constant \(c_0\) and whole zero-cap term
\(-mn_0\log t/(2\pi i)\) are in \(U_0\). Put
\[
 \tau(v)=\left(\frac{\varepsilon\Gamma}{b},
 \frac{U_0+\varepsilon C-\varepsilon Mv}{b}\right),
 \quad
 \sigma(v)=\left(\varepsilon\Gamma v,
 v(U_0+\varepsilon C)-\frac{\varepsilon M}{2}(v^2-v)\right).
\tag{JFP.36}
\]
Then \(\sigma(v+1/b)-\sigma(v)=\tau(v)+
\varepsilon M(b-1)\delta/(2b^2)\). The coordinate
\(\xi=z-\sigma(v)\) closes for \(cE_r\). The entire
intermediate products
\(C_s(v)=\exp(\sigma(v)-\sigma(v+s))\) and
\(P_j(v)=\exp(\sigma(v)-\sigma(v+j/b))\) satisfy
\(C_s(v+j/b)P_j(v)=C_{j/b+s}(v)\).
Their negative exponents are the inverse. At \(j=b\) this
is the actual inverse one-gamma closing, with every derivative
matched on the end collars.

For a basis \(P\), retain all four full coefficient units in
the conjugated period matrix. If its two multiplicative factors
are \(A,D\), the actual vertex and Cox arrows are
\[
\begin{split}
 (a,b_1,c_1)&\mapsto((AD)^{-1}a,Ab_1,Dc_1),\\
 (u,v_1,w_1)&\mapsto(ADu,D^{-1}v_1,A^{-1}w_1),\\
 (a_0,a_1,c_0,c_1)&\mapsto(a_0,Aa_1,(AD)^{-1}c_0,Dc_1).
\end{split}
\tag{JFP.37}
\]
They preserve \(q\), the Cox relation and both ruling charts;
reciprocal factors give the inverse on the exceptional curve
as well. Substitution into the three complete transition rows
\((\kappa\alpha b_1c_1,\alpha^{-1}ab_1,\kappa^{-1}b_1^{-1})\),
\((g\eta b_1c_1,\eta^{-1}c_1^{-1},g^{-1}ac_1)\), and
\((a^{-1},ab_1,ac_1)\) gives the same factor on both sides.
For \(F_r\) these factors are
\(r^{N(P^{-1})_{12}},r^{N(P^{-1})_{22}}\). Every full
\(q\)-dependent unit and higher normal jet is retained by the
normal rescaling; its cutoff flow fixes the protected collars.
The ordinary and cap maps are exactly JFP.3--5. Thus (JFP.29)
is an isomorphism of the entire parameter families and every
intermediate seam.

The free parameter action has quotient \(w=r^b\),
\(\lambda=w^{J_0}\). Its complete \(r\)-family is the literal
degree-\(b\) pullback. The physical line, cap identification and
stable tangent pull back by the maps just written. Integration
of their integral characteristic classes therefore multiplies
by \(b\). This gives the primitive integral comparison before
any reduction modulo 28.

\par

\medskip\noindent\textbf{JFP.9. Page, physical pair and native product}

\par\nobreak

The actual original canonical page over ordinary binding
\(\tau\) is
\[
 \eta^h-\tau\xi^m=1,\qquad
 t=\tau+\xi^{-m}=\eta^h\xi^{-m}.
\tag{JFP.38}
\]
The reciprocal roots are \(\xi^m=(t-\tau)^{-1}\),
\(\eta^h=t/(t-\tau)\), with the original cap labels.
The degree is \(M\), with branch indices \(h,m,M\) at
\(0,\infty,\tau\); its genus is \((m-1)(h-1)/2\).
Pull back the original source by these maps and apply the
actual root resolutions in JFP.1. The relative Picard
evaluation compares its pulled-back point line with the
translated degree-zero line on each fibre. Its component
degrees vanish by (JFP.6--10), so the remaining evaluation
divisor is an integral sum of whole fibres. Division by their
base equations gives an invertible map at every singular
component. Rigidification at O fixes its scalar. The residual
base line changes the class by \(zF\), with \(NF=F^2=0\).
Its boundary value is the specified one and extends over the
ordinary binding disc; the full relative gauge has zero square
ambiguity. The tangent comparison uses the same source map.

The original section height is \(1/6\). The I1 has \(M\)
preimages, infinity has \(h\) root fibres, and zero has \(m\).
Thus
\[
 H=Ma^2/6,\quad N^2=-H,\quad
 \chi=(M+h e_\infty+m e_0)/12,
 \quad(e_\infty,e_0)=(0,0),(0,3),(4,0),(4,3).
\tag{JFP.39}
\]
The full signed binding graph is
\(q_b=(t-\tau)(t/\tau)^{-\varepsilon mn_0}\), whose
derivative at \(\tau\) is one. Retain its constant phase
and all unit derivatives. The normal angles obey
\(\alpha_{\rm normal}=\phi-M\beta\), \(\beta=\varepsilon\gamma\),
and the source graph has \(\beta=\theta+g(p)\).
The K30 triangular source/tangent shear subtracts the full
derivatives of \(g\) and that binding unit. Equivalently the
whole origin difference is
\(\varepsilon\theta\gamma-\varepsilon M\theta(\theta-1)\delta/2\).
Its positive parameter orientation is retained for both signs.
The actual K68 dual-line and inverse-clutch Hopf map consequently
gives the primitive pair and its held-quaternion native product
\[
 (k,\ell)=(3H/b,24\chi/b),\qquad
 (\mu,\nu)=(6\chi/b,3H/b-6\chi/b).
\tag{JFP.40}
\]
The cover in JFP.8 proves this as an equality of integral
characteristic pairs. The native left/right relation is
\(k=\mu+\nu,\ell=4\mu\) under that actual product map.
The four primitive indices are
\(0,kl,k(4l+2),35kl+22k+8l+5\). In full pairs they are
\((kl,4kl)\), \((2l(4k+1),8l(k+1))\),
\((9k(3l+2),12k(l+2))\), and
\((18(3l+2)(4k+1),24(kl+2k+l+1))\).

\par

\medskip\noindent\textbf{JFP.10. Actual fan maps and independent full target systems}

\par\nobreak

Use the literal K71 identical-fan maps (column swap and the
six right triangle rotations) on the absolute semigroups.
They are inverse holomorphic maps, identity on the original
source/Cartier/cap overlaps. Euclidean subtraction of the
positive pair terminates at one vertical \(T_j\). Reversing
uses the actual two flops
\(Q_0=\left(\begin{smallmatrix}1&-1\\0&1\end{smallmatrix}\right)\),
\(Q_1=\left(\begin{smallmatrix}1&0\\-1&1\end{smallmatrix}\right)\).
Their dense characters are \((ZV,V)\), \((Z,ZV)\), their
first triples \((a/c,bc,c)\), \((a/b,b,bc)\), and their
two ruling charts
\((n_0/n_1,zn_1,n_1)\), \((n_1/n_0,zn_0,n_0)\).
The common rank-one matrix and \(q=zn_0n_1\) are retained.
Every full coefficient is transformed by
\(A'_{ik}=\prod_{hl}A_{hl}^{(Q^{-1})_{ih}Q_{lk}}\), with
inverse obtained from \(Q^{-1}\). The source point line and
two original caps therefore carry through each finite word.

For the scalar family of primitive winding \(J=J_0\), the
complete ordinary product has determinant winding \(2J\).
Its native quaternion factors are \(H_J QH_{L_i-J}\), with
the periodic coefficient logarithms and all normal jets retained.
At the second vertex the full ordinary height is
\(qA_{11}(q)A_{22}(q)\); its inverse first recovers the original
\(q\) and divides at that recovered value. The cutoff unit
product similarly has height \(q\exp[\rho(k_{i0}+k_{i1})(q)]\),
whose uniform inverse follows from its unit derivative on a
smaller protected disc. The held Hopf map is
\((UH_JU^{-1})q_HH_{L_i-J}\), of left degree \(2J\) and
right degree zero. Its inverse cylinder is \(L^{-2J}\).
The full circle and derivative prisms therefore have exactly
the K73 data below, with their base and normal Spin endpoints
paired. This is the ordinary-to-native identification carried
by (JFP.37), including the full source-dependent coefficients.

Set \(A=S_0=H/(2b)-\chi/b\), \(\mu=6\chi/b\).
At \(T_j\), \(\nu_j=6A-3Jj(j-1)\),
\[
\begin{split}
 V&=\begin{pmatrix}
 2J-2\mu-2\nu_j&2J(1-2j)&2J(2j-1)\\
 2J(1-2j)&2J&0\\2J(2j-1)&0&2J
 \end{pmatrix},\\
 D&=(2J-2\mu+2\nu_j,2J(2j-1),2J(1-2j)),
 \qquad S=\nu_j/6.
\end{split}
\tag{JFP.41}
\]
At a directed edge select \(a_*=V_{ii},b_*=D_i\) and
\(P'=PQ_i\). Put
\[
\begin{split}
 B&=2J(P')^{-1}e_2e_2^tP',\quad
 L=(B_{22}-B_{21},B_{11}-B_{12},B_{22}),\\
 (x,y)^t&=(P')^{-1}e_2,\quad
 w=(x,y,-x-y),\quad d=(-x-2y,-2x-y,x-y),\\
 I_r&=Jw_r+d_r(L_r-J),\quad W_r=Jw_r-d_r(L_r-J).
\end{split}
\tag{JFP.42}
\]
The actual common-link map conjugates the first quaternion by
the held second quaternion. The selected native reset gives
\(V'_{22}=a_*\), \(D'_2=2a_*-b_*\),
\(S'=S+(b_*-a_*)/4\). Together with all equations
\[
 V'(w_ie_i-w_0e_0)=I_ie_i-I_0e_0,\quad
 w_iD'_i-w_0D'_0=W_i-W_0,\qquad i=1,2,
\tag{JFP.43}
\]
these determine all six symmetric entries and all three
derivatives. The independent executable solves these complete
systems by Fraction row elimination, without importing the
primary or another fan program. Each target has rank nine,
and every redundant original equation is checked. The resulting
entries agree with
\(V'_{rs}=(I_r/w_r)\delta_{rs}+C/(w_rw_s)\),
\(D'_r=(W_r+E)/w_r\), where
\(C=w_2^2a_*-w_2I_2\), \(E=w_2(2a_*-b_*)-W_2\).
All target weights are nonzero in the finite canonical word.
The inverse reset gives the negative of its forward right
increment. Thus these are the actual all-fan returns rather
than only solutions of a numerical recursion.

\par

\medskip\noindent\textbf{JFP.11. Exact unbounded images and all scalar paths}

\par\nobreak

Linearity of the complete target systems in \((J,\mu,A)\)
gives the divisibility upper bounds from actual admissible
unscaled profiles. For \(a=1\), all inputs are \(kl\) times
the \((k,l)=(1,1)\) source. The two full selected pairs along
\(Q_0Q_0\) are \((2,2)\), then \((-4,8)\), and the
return is \(3kl\). Since 3 is a unit modulo 28, the image is
\(kl\mathbb Z/28\).

For \(a=2\), all inputs are \(l\) times the actual \((k,1)\)
profile. Its vertical coefficients are
\(S_j/l=k-(4k+1)j(j-1)/2\). For odd \(k\), choose
\(j=0\) when \(k\ne0\pmod7\), and \(j=4\) otherwise;
the coefficients are \(k\), \(-23k-6\). For even \(k\),
choose \(j=2\) when \(k\ne2\pmod7\), and \(j=3\)
otherwise; the coefficients are \(-3k-1\), \(-11k-3\).
Each is odd and nonzero modulo 7. The image is
\(l\mathbb Z/28\).

For \(a=3\), admissibility forces \(l\) odd. All inputs
are \(k\) times the actual \((1,l)\) profile. The independent
three full target solves along \(Q_0^3\), applied to
\(\mu/k=3(l+2)\), \(A/k=4l+2\), give
\[
 S(Q_0^3)/k=15\cdot3(l+2)+127(4l+2)
 =553l+344=7(79l+49)+1.
\tag{JFP.44}
\]
It is odd and nonzero modulo 7. The image is
\(k\mathbb Z/28\).

For \(a=6\), write \(A=35kl+22k+8l+5\) and
\(J=(3l+2)(4k+1)\). When \(kl\) is even, the two odd
integers \(A,A-6J\) have a common zero modulo 7 only at
\((l,k)=(4,5)\pmod7\): substitution yields
\(2(l-4)^2=0\). When \(kl\) is odd, use the odd integers
\(A-J,A-3J\); their common zero again forces that pair,
by \(3l^2+11l+6=3(l-4)^2\pmod7\).
It makes both \(m,h\) divisible by 7, contradicting their
coprimality. Hence one displayed vertical fan has a return
unit modulo 28, giving the whole group. These finite residue
identities are polynomial identities in the unrestricted
orders, rather than bounded tests of the orders.

The complete table is
\[
\begin{array}{c|c|c}
(a,b)&\text{all holomorphic winding values}&
\text{all-fan, all-loop return image}\\\hline
(1,6)&2kl\mathbb Z&kl\mathbb Z/28\\
(2,3)&l(4k+1)\mathbb Z&l\mathbb Z/28\\
(3,2)&2k(3l+2)\mathbb Z&k\mathbb Z/28\\
(6,1)&(3l+2)(4k+1)\mathbb Z&\mathbb Z/28.
\end{array}
\tag{JFP.45}
\]
For a full holomorphic scalar path, its logarithm on the
universal strip obeys \(L(\theta+1)-L(\theta)=2\pi iJ\).
Interpolate to \(L(0)+2\pi iJ\theta\). Convexity of the
logarithmic modulus bound keeps it in the original contraction
disc. Substitute this full logarithm into (JFP.20--25), retaining
the original finite characters, integer periods and references,
and use (JFP.36--37) on every selected chart. Common strict
contraction tubes exist on the compact interpolation interval.
Local parameter lifts, patched relative to the specified end
collars, integrate to a proper relative isotopy. The endpoint
ratio is a physical scalar by JFP.7; multiplication by
\(\exp(u\log c)\) gives its extending isotopy. Winding
\(sJ_0\) is attained by the literal annular pullback
\(w\mapsto w_*(w/w_*)^s\), with inverse maps for negative
\(s\). At a fixed fan the complete return set is therefore
\(sS(P)\bmod28\), and (JFP.45) follows for all paths and
all holomorphic closings.

\par

\medskip\noindent\textbf{JFP.12. Mapping-torus orientations and audit scope}

\par\nobreak

The full disk-transfer and seam maps retained in K114--115 use
the physical inverse-clutch convention
\(\lambda_{\rm physical}(F)=-\Sigma_{F^{-1}}=\Sigma_F\).
For \((x,1)\sim(\Phi(x),0)\), the oriented mapping torus
is \((S^1\times S^6)\#\Sigma_{sS(P)}\). Reversing the
seam uses the actual inverse endpoint and gives
\(\Sigma_{-sS(P)}\). For each of the two spin structures
the corresponding Eells--Kuiper values are
\(sS(P)/28\) and \(-sS(P)/28\) in \(\mathbb Q/\mathbb Z\).
These signs are those of the binding graph and Hopf map in
JFP.9; no change of the signed sphere equation changes the
chosen positive scalar-parameter orientation.

The independent certificate includes the four symbolic root
and page rows, all signed norm and local character identities,
the whole finite good-root Miller functions, the actual
reduced-IV marked graphs, the integral source lattice, the
bounded-origin and vertex squares, and 867 fresh complete
nine-variable systems. The first executed anchor-only result
and its script are preserved before this full-source refresh.
The complete primary RJC.1--13 and all named stable anchors
have been read. The a=3 graph constant correction is retained
in (JFP.10); no further geometric discrepancy was found.
The full source/cap/reference and negative-lift arguments,
both mapping-torus signs, and the unbounded image proof are
part of this peer, beyond the finite arithmetic checks.

The scope is precisely (JFP.1--3), the original two additive
caps, zero additional ordinary total norm, every retained
finite character and annular reference, and all selected
lattice-basis fans. Other source-moduli paths and redistributed
nonzero ordinary clutch totals have their own full comparison
problems and are not asserted by this proof.

% End expanded verified cp36 input



\par\medskip\noindent\textbf{Exact construction and audit dependencies.}\par
The complete original F46 joint-cap construction is
\path{notes/part3_iv_star_iii_joint_working.md}.
Its actual point/O graph, normalized blowups, primitive cap generators
and integral sphere maps are retained in both supporting proofs.
The K67 shorter-parameter family is
\path{notes/part3_shorter_scalar_vertical_fans_working.md};
the exact K73 all-fan product and prism construction is
\path{notes/part3_scalar_all_fans.tex}.
The completed K114--115 source and endpoint comparisons are
\path{notes/part3_e6_iii_endpoint_fourfold_peer.md} and
\path{notes/part3_e6_complete_scalar_returns_working.md}.
The canonical two-cap page and its signed source graph are
\path{notes/part3_joint_additive_returns_working.md}.

The primary executable is
\path{code/part3_root/e6_iii_joint_endpoint_derivation.py}.
The separate independent executable is
\path{code/part3_fourfold/e6_iii_joint_endpoint_peer.py}.
The independent source and endpoint derivation, fresh full target
systems, original-coordinate marked graphs and complete residue
witnesses are locked with their entire supporting proofs.
Every source, executable, exact certificate and full geometric audit
dependency is listed in
\path{records/k118_export_manifest.json}.
The bounded executions support the displayed morphisms; the source
extension, finite quotient, terminating fan maps and exact modular
identities prove the full stated domains.


% End expanded verified cp36 input

\end{proof}


\subsection{OS45 ordinary-cap collision meridians with the full original reference}
